\chapter{West Chadic (29/84)}

West Chadic languages are divided into two subbranches. Descriptions are available for the grammar of 22 of the 46 West Chadic A languages. Descriptive coverage of the West Chadic B languages is much poorer with descriptions available for only seven of 38 languages. 

\section{West Chadic A (22/46)}

The 46 West Chadic A languages are divided into four groups. In each group, grammatical descriptions are available for half of the languages or more, with the exception of West Chadic A3. 

\subsection{West Chadic A1 (1/2)}

No description of the morphosyntax of Gwandara (gwan1268) is available, so only Hausa is included here.

\paragraph*{Hausa (haus1257)} \label{haus}

Information on Hausa grammar comes mainly from a reference grammar \citep{newmanHausaLanguageEncyclopedic2000}. 

\largerpage
\subparagraph{Directional verbs}~
~


\begin{table}[H]
	\caption{Hausa directional verbs} 
	\label{tab:hausV}
	\begin{tabular}{lll} % add l for every additional column or remove as necessary
		\lsptoprule
		DIR & verb & gloss   \\ %table header
		\midrule
		ITIVE & \textit{tafi} & `go' \\
		VENT & \textit{zo} & `come' \\
		UP & \textit{hawa} & `go up, climb' \\
		DOWN & \textit{sauka} & `descend' \\
		INTO& \textit{shiga} & `enter' \\
		OUT & \textit{fita} & `exit' \\
		\lspbottomrule
	\end{tabular}
\end{table}

\subparagraph{Directional extensions}~



\begin{table}[H]
	\caption{Hausa directional morphology}
	\label{tab:haus}
	\begin{tabular}{llll} % add l for every additional column or remove as necessary
		\lsptoprule
		forms & source gloss & directional  & other functions  \\ %table header
		\midrule
		\textit{-ō}  & Grade 6  &   \textsc{vent} &  \textsc{subs.vent, loc}  \\
		\lspbottomrule
	\end{tabular}
\end{table} 

The Hausa ventive suffix is a long vowel \textit{-ō} traditionally called ``Grade 6'' in the study of Hausa verbal morphology. ``The ventive ending generally denotes action or movement in the direction of the speaker (or any other pragmatically established deictic center)…'' \citep[663]{newmanHausaLanguageEncyclopedic2000}, as in (\ref{ex2:haus1}) and (\ref{ex2:haus2}).

\ea\label{ex2:haus1}
\langinfo{Hausa verb \textit{fit} `exit' with ventive suffix}{}{\cite[663]{newmanHausaLanguageEncyclopedic2000}} \\
\gll fit-ō \\
exit-\gl{vent} \\
\glt `come out'

\ex\label{ex2:haus2}
\langinfo{Hausa verb \textit{jany} `drag' with ventive suffix}{}{\cite[663]{newmanHausaLanguageEncyclopedic2000}} \\
\gll jany-ō \\
drag-\gl{vent} \\
\glt `drag in this direction'
\z 

The same suffix also ``sometimes indicates `do some action and come'{''} \citep[663]{newmanHausaLanguageEncyclopedic2000}, in other words, subsequent associated motion, as in (\ref{ex2:haus3}) and (\ref{ex2:haus4}). 

\ea\label{ex2:haus3}
\langinfo{Hausa verb \textit{shāf} `whitewash' with ventive suffix}{}{\cite[663]{newmanHausaLanguageEncyclopedic2000}} \\
\gll nā	shāf-ō	bangō	 \\
I	whitewash-\gl{vent}	wall \\
\glt `I whitewashed the wall and came back.'

\ex\label{ex2:haus4}
\langinfo{Hausa verb \textit{said} `sell' with ventive suffix}{}{\cite[664]{newmanHausaLanguageEncyclopedic2000}} \\
\gll said-ō \\ 
sell-\gl{vent} \\
\glt `sell and come back'
\z 

In one example, shown in (\ref{ex2:haus5}), the interpretation of the ventive suffix is subsequent caused accompanied motion.

\ea\label{ex2:haus5}
\langinfo{Hausa verb \textit{say} `buy' with ventive suffix}{}{\cite[663]{newmanHausaLanguageEncyclopedic2000}} \\
\gll yā	say-ō	nāmā̀ \\
he	buy-\gl{vent}	meat \\
\glt `He bought some meat and brought it back here.'
\z 

In at least one case, a verb with a ventive suffix can be interpreted as either directional or subsequent associated motion. In (\ref{ex2:haus6}), the verb `put' with a ventive suffix can either be directional in a transitive sense where the object is moving towards the deictic center, or as subsequent associated motion where the verb is interpreted (intransitively) as `getting dressed'. 

\ea\label{ex2:haus6}
\langinfo{Hausa verb \textit{sāw} `put' with ventive suffix}{}{\cite[662]{newmanHausaLanguageEncyclopedic2000}} \\
\gll sāw-ō \\
put-\gl{vent} \\
\glt `put here, put on and come'
\z 

The are also some uses of ventive forms that have no motion meaning including \textit{ɓull-ō} `appear suddenly' and \textit{far̃faɗ-ō} `revive' \citep[663]{newmanHausaLanguageEncyclopedic2000}.

There is another verbal suffix that is sometimes translated with a directional meaning. The suffix \textit{-ar̃} is the ``Grade 5'' form which is traditionally labeled a causative form. \citet{jaggarHausaGrade52014} discusses the many uses of this suffix, including a few examples where the suffix is used with a transitive verb and given an itive translation such as \textit{jeef-ar̃} `throw away', \textit{zub-ar̃} `pour/throw away', \textit{tuur-ar̃} `push away' and \textit{wur̃gar̃} `fling away, dismiss' (from \textit{wur̃gàa} `throw'). In other cases, the same suffix has an outward directional translation such as \textit{gus-ar̃} `move out, remove' and \textit{juuy-ar̃} `dump out'.

\citet{parsonsFurtherObservationsCausative1962} observes that these forms carry a sense of ``disposal'' or ``riddance''. It is therefore unclear whether these forms express literal itive direction, or if the itive forms in the English translations are metaphorical ways that English speakers refer to acts of disposal. Since there are no unambiguous cases of a literal motion use of this suffix, Hausa is not considered to have an itive or outward directional extension.

\subparagraph{Caused accompanied motion (CAM)}



The Hausa verb \textit{kai} `take' has a caused accompanied motion interpretation with a ventive suffix, as in (\ref{ex2:haus7}). 

\ea\label{ex2:haus7}
\langinfo{Hausa verb \textit{kāw} `take' with ventive suffix}{}{\cite[663]{newmanHausaLanguageEncyclopedic2000}} \\
\gll kāw-ō \\
take-\gl{vent} \\
\glt `bring'
\z 

A similar interpretation is found with the ventive form of the verb `take out', \textit{fidd-ō} `bring out' \citep[664]{newmanHausaLanguageEncyclopedic2000}. Note in (\ref{ex2:haus5}) that the interpretation also includes caused accompanied motion (i.e., brought the meat), although this is an implicature rather than an entailment. 

\subparagraph{Buy-sell verbs}



	The verb \textit{sàyaa} `buy' in a causative form (``Grade 5'') becomes \textit{sayar̃} `sell' \citep[][]{jaggarHausaGrade52014}. The ventive form of the verb remains `buy' as in (\ref{ex2:haus5}).

\subsection{West Chadic A2 (14/22)}

There are two subgroups of West Chadic A2: Boleic and Tangalic.

\subsubsection{Boleic, West Chadic A2 (10/15)}


Relevant publications are available for ten of the Boleic languages excluding Daza (daza1244),
Deno (deno1239),
Kubi (kubi1239),
Kholok (khol1240)
and Maaka (maak\-1236). Information about many Boleic languages is primarily from notes in \citet{schuhBoleTangaleLanguagesBauchi1978}, based on a few hours of elicitation with speakers of each language.

\paragraph*{Karekare (kare1348)} \label{kare} 

Information on the grammar of Karekare is from \citet{zochVerbalmorphologieBoleTangaleSprachenWesttschadisch2014}, unpublished manuscripts of Russell Schuh \citep{schuhKarekareVerbalParadigmsn.d.,schuhKarekareVerbalSystemn.d.}, a recent MA thesis \citep{abareVerbalMorphosyntaxKaraiKarai2020}, and a dictionary \citep{gamboKarekareEnglishHausaDictionary2004}. 

\subparagraph{Directional verbs}~



\begin{table}[H]
	\caption{Karekare directional verbs}
	\label{tab:kareV}
	\begin{tabular}{lll} % add l for every additional column or remove as necessary
		\lsptoprule
		DIR & verb & gloss   \\ %table header
		\midrule
		MOTION & \textit{ndu} & `go, come' \\
		ITIVE & \textit{wā̀lu} & `go' \\
		INTO& \textit{rā} & `enter' \\
		OUT & \textit{fàtā} & `exit' \\
		DOWN & \textit{(w)ùru} & `descend' \\
		\lspbottomrule
	\end{tabular}
\end{table}

Verb forms in Table \ref{tab:kareV} are from \citet{gamboKarekareEnglishHausaDictionary2004}.

\subparagraph{Directional extensions}~



\begin{table}[H]
	\caption{Karekare directional morphology}
	\label{tab:kare}
	\begin{tabular}{llll} % add l for every additional column or remove as necessary
		\lsptoprule
		forms & source gloss & direction  & other functions  \\ %table header
		\midrule
		\textit{-(n)ee, -tu}  &    &   \textsc{vent} &   ??  \\
		\lspbottomrule
	\end{tabular}
\end{table} 

The available descriptions of Karekare verbal morphosyntax are somewhat limited. \citet[viii]{gamboKarekareEnglishHausaDictionary2004} and \citet[287]{zochVerbalmorphologieBoleTangaleSprachenWesttschadisch2014} describe the ventive as having the form \textit{-(n)ee} in Perfective aspect and \textit{-tu} in Imperfective aspect and Subjunctive mood. Schuh's manuscripts and \citet[48]{abareVerbalMorphosyntaxKaraiKarai2020} confirm these forms in several paradigms without providing translations or example sentences.

The only glossed and translated linguistic examples available are from the field notes of Suleiman Abare.\footnote{Thanks to Matthew Harley for making these data accessible.} These include a clear distinction between a motion verb with and without a ventive suffix. The verbs in (\ref{ex2:kare1}) and (\ref{ex2:kare3}), without a ventive suffix, have an itive interpretation. In contrast, those same verbs in (\ref{ex2:kare2}) and (\ref{ex2:kare4}) with a ventive suffix are given a ventive interpretation---indicating that the direction of the motion is toward the deictic center.

\ea \label{ex2:kare1}
\langinfo{Karekare verb \textit{r} `enter' without ventive suffix}{}{\cite[]{abareTCNN16Suleiman2016}} \\
\gll lewi ra-kau a benu	\\
child enter-\gl{pfv} \gl{prep} room	\\
\glt `The child went into the room.'

\ex\label{ex2:kare2}
\langinfo{Karekare verb \textit{r} `enter' with ventive suffix}{}{\cite[]{abareTCNN16Suleiman2016}} \\
\gll lewi re-ne-kau a benu \\
child enter-\gl{vent}-\gl{pfv} \gl{prep} room	\\
\glt `The child came into the room.'

\ex\label{ex2:kare3}
\langinfo{Karekare verb \textit{yar} `move' without ventive suffix}{}{\cite[]{abareTCNN16Suleiman2016}} \\
\gll raɓjau yi yar-kau ka gi-nau \\
lady the move-\gl{pfv} from place-my	\\	
\glt `The lady moved away from me.'

\ex\label{ex2:kare4}
\langinfo{Karekare verb \textit{yar} `move' with ventive suffix}{}{\cite[]{abareTCNN16Suleiman2016}} \\
\gll raɓjau yi yar-ne-kau a gi-nau	\\
lady the move-\gl{vent}-\gl{pfv} \gl{prep} place-my	\\
\glt `The lady moved towards me.'
\z 

The verb root \textit{nd} is a general motion verb root that does not indicate direction, but can be modified by the ventive suffix to indicate ventive direction, as in (\ref{ex2:kare5}) and (\ref{ex2:kare6}). There is some variability in the form of the vowel at the end of the verb root. Before a ventive suffix it is often \textit{e}, but elsewhere \textit{a} appears to be more common, and this vowel is retained in (\ref{ex2:kare6}).  

\ea\label{ex2:kare5}
\langinfo{Karekare verb \textit{nd} `go/come' without ventive suffix}{}{\cite[]{abareTCNN16Suleiman2016}} \\
\gll badine yi nde-ne-kau	\\
	 girl the go/come-\gl{vent}-\gl{pfv}	\\
\glt `The girl has come.'

\ex\label{ex2:kare6}
\langinfo{Karekare verb \textit{nd} `go/come' with ventive suffix}{}{\cite[]{abareTCNN16Suleiman2016}} \\
\gll musa ka dauda nda-ne-kau \\
Moses and David go/come-\gl{vent}-\gl{pfv}	\\
\glt `Moses and David came.'
\z 

\subparagraph{Caused accompanied motion (CAM)}



A pair of example sentences show that the ventive suffix can be used with the verb \textit{wula} `carry' to indicate ventive caused accompanied motion. Without the ventive suffix, in (\ref{ex2:kare7}), the translation is itive. With the ventive suffix, in (\ref{ex2:kare8}), the translation is ventive.

\ea\label{ex2:kare7}
\langinfo{Karekare verb \textit{wula} `carry' without ventive suffix}{}{\cite[]{abareTCNN16Suleiman2016}} \\
\gll mendau yi 'wula waɗa a mizitau	\\ 
woman the carry food \gl{prep} male \\
\glt `The woman carried food to her husband.'

\ex\label{ex2:kare8}
\langinfo{Karekare verb \textit{wula} `carry' with ventive suffix}{}{\cite[]{abareTCNN16Suleiman2016}} \\
\gll mendau yi 'wule-ne waɗa a mizitau \\
woman the carry-\gl{vent} food \gl{prep} male	\\
\glt `The woman brought food to her husband.'
\z 

\subparagraph{Buy-sell verbs}



There are two verb roots in Karekare: \textit{jànā} `buy' and \textit{bàutu} `sell' \citep[6, 28, 62, 81]{gamboKarekareEnglishHausaDictionary2004}. 

\paragraph*{Beele (beel1236)}  \label{beel}

A very limited amount of information about the grammar of Beele is from a short grammar sketch \citep{schuhBoleTangaleLanguagesBauchi1978}. 

\subparagraph{Directional verbs}~



\begin{table}[H]
	\caption{Beele directional verbs} 
	\label{tab:beelV}
	\begin{tabular}{lll} % add l for every additional column or remove as necessary
		\lsptoprule
		DIR & verb & gloss   \\ %table header
		\midrule
		ITIVE & \textit{jíi-kò} & `go' \\
		VENT & \textit{ndûŋ} & `come' \\
		OUT & \textit{fétí-kò}  &  `go out' \\
		\lspbottomrule
	\end{tabular}
\end{table}

\subparagraph{Directional extensions}


%\begin{table}[H]
%	\caption{Beele directional morphology}
%	\label{tab:beel}
%	\begin{tabular}{llll} % add l for every additional column or remove as necessary
%		\lsptoprule
%		forms & source gloss & direction  & other functions \\ %table header
%		\midrule
%		\textit{-u}  &  ventive  &   ?? &   ??  \\
%		\lspbottomrule
%	\end{tabular}
%\end{table} 
On Beele, \citet[21]{schuhBoleTangaleLanguagesBauchi1978} writes: ``The ventive has the form Verb Root + u''. The only examples given are \textit{pán-ũ̂} `carry-\textsc{vent}' and \textit{bàah-ũ̂} `shoot-\textsc{vent}'. Presumably both are cases of ventive directional meaning, but no translations are given. No further relevant data on Beele is available \citep[cf.][292]{zochVerbalmorphologieBoleTangaleSprachenWesttschadisch2014}. Due to this lack of examples, Beele is not counted as a language with a ventive directional extension.

\subparagraph{Caused accompanied motion (CAM)}



\citet[26]{schuhBoleTangaleLanguagesBauchi1978} includes a lexicon with the verbs \textit{àlũ̂} `bring' (presumably ventive?) and \textit{ál-kò} `take (from one place to another)' but no other information is given about directed caused accompanied motion.

\subparagraph{Buy-sell verbs}



The Beele wordlist includes \textit{gòjjú-kò} `buy' but no verb for `sell'. 
 
\paragraph*{Bole (nucl1695)} \label{bole} 

Bole is described in a dissertation by \citet{gimbaBoleVerbMorphology2000} and an unpublished draft of a grammar \citep{schuhBoleReferenceGrammar2005}. There is also a Bole dictionary \citep{gimbaBoleenglishhausaDictionaryAlhaji2004}.

\subparagraph{Directional verbs}~



\begin{table}[H]
	\caption{Bole directional verbs} 
	\label{tab:boleV}
	\begin{tabular}{lll} % add l for every additional column or remove as necessary
		\lsptoprule
		DIR & verb & gloss   \\ %table header
		\midrule
		MOTION & \textit{ndà} & `go/come' \\
		DOWN & \textit{yàwwu} & `descend, go downward' \\
		INTO& \textit{rá} & `enter' \\
		OUT & \textit{pàtà} & `go out' \\
		\lspbottomrule
	\end{tabular}
\end{table}

\subparagraph{Directional extensions}~



\begin{table}[H]
	\caption{Bole directional morphology}
	\label{tab:bole}
	\begin{tabular}{llll} % add l for every additional column or remove as necessary
		\lsptoprule
		forms & source gloss & direction  & other functions  \\ %table header
		\midrule
		\textit{-n, -tV, -aakoo}  &  \textsc{vent}  &   \textsc{vent} &  \textsc{subs.vent}, \textsc{loc.here}, \textsc{ben}  \\
		\lspbottomrule
	\end{tabular}
\end{table} 

\citet{gimbaBoleVerbMorphology2000} and \citet{schuhBoleReferenceGrammar2005} describe ventive verbal forms in Bole. There are three ventive suffixes distributed according to the tense-aspect of the verb: \textit{-n} with Perfective verbs, \textit{-tV} (where \textit{V} stands for a vowel) with Subjunctive{\slash}Imperative verbs and \textit{-aakoo} with Future and Progressive verbs. The Future and Progressive/Habitual are said to be ``nominal in origin'' \citep[307]{schuhChadicCornucopia2017}. Nominal verbs are segmentally identical, but the verb is preceded by a marker \textit{á}, and they differ in tone.

As in other Chadic languages, the ventive form with motion verbs can give a ventive directional meaning, as in (\ref{ex2:bole3}) and (\ref{ex2:bole4}).

\ea\label{ex2:bole3}
\langinfo{Bole verb \textit{pàt} `exit' with ventive suffix}{}{\cite[135]{gimbaBoleVerbMorphology2000}} \\
\gll íshí	à	pàt-àakóo \\
he	??	exit-\gl{vent} \\
\glt `He will come out.'

\ex\label{ex2:bole4}
\langinfo{Bole verb \textit{tùbbú} `push' with ventive suffix}{}{\cite[135]{gimbaBoleVerbMorphology2000}} \\
\gll ítá	tùbbú-n	kúlà \\
she	push-\gl{vent}	calabash \\
\glt `She pushed the calabash here.'
\z 

With non-motion verbs, the ventive suffix can have a ventive subsequent motion interpretation, as in (\ref{ex2:bole5}). 

\ea\label{ex2:bole5}
\langinfo{Bole verb \textit{gàtt} `get tired' with ventive suffix}{}{\cite[139]{gimbaBoleVerbMorphology2000}} \\
\gll à gàtt-àakóo \\
?? get.tired-\gl{vent} \\
\glt `He will get tired and come.'
\z 

\citet{gimbaBoleVerbMorphology2000} and \citet{schuhBoleReferenceGrammar2005} also includes examples translated with subsequent caused accompanied motion, i.e., `bring', as in (\ref{ex2:bole6}) and (\ref{ex2:bole1}). 

\ea\label{ex2:bole6}
\langinfo{Bole verb \textit{bèsé} `shoot' with ventive suffix}{}{\cite[135]{gimbaBoleVerbMorphology2000}} \\
\gll màté	bèsé-t-tùu	sá \\
they	shoot-\gl{vent}-\gl{vent}	\gl{neg} \\
\glt `They should not shoot (it) and bring.'
\z

The caused accompanied motion element of the translation appears to be an implicature rather than part of the meaning of the ventive suffix in these contexts. Compare the ventive form of ‘tie’ in (\ref{ex2:bole1}) with a ventive subsequent motion translation to the same verb in (\ref{ex2:bole2}) with a locative translation not explicitly indicating any motion. The translations of many examples of ventive forms in Bole do not include any motion meaning, but rather a locative meaning translated by the adverbial `here'. This may indicate some variation in the use of these markers, or some translations might simply be less explicit. 

\ea\label{ex2:bole1}
\langinfo{Bole verb \textit{ngòrú} `tie' with ventive suffix}{}{\cite[143]{gimbaBoleVerbMorphology2000}} \\
\gll íshí	ngòrú-t-tú	dóoshó \\
he	tie-\gl{vent}-\gl{vent}	horse \\
\glt `That he tie a horse and bring here.'

\ex\label{ex2:bole2}
\langinfo{Bole verb \textit{ngòrú} `tie' with ventive suffix}{}{\cite[142]{gimbaBoleVerbMorphology2000}} \\
\gll íshí	ngòrú-t-tùu-yì \\
he	tie-\gl{vent}-\gl{vent}-\gl{obj} \\
\glt `That he tie (it) here.'
\z 

Gimba points out a particular semantic pattern when the ventive is used with a pronominal suffix. Pronominal suffixes can either have a direct or indirect object function in Bole. When a ventive co-occurs with a pronominal suffix, the interpretation is that the referent of the pronominal is a beneficiary. In most cases, the translations Gimba provides still include a ventive meaning alongside the beneficiary, as in (\ref{ex2:bole7}) and (\ref{ex2:bole8}) (cf. \ref{ex2:bole1} and \ref{ex2:bole2}).\footnote{Note that the suffix \textit{-yi} glossed \textsc{obj} in (\ref{ex2:bole2}) is not considered a pronominal suffix.}

\ea\label{ex2:bole7}
\langinfo{Bole verb \textit{ngòrú} `tie' with ventive suffix and pronominal suffix}{}{\cite[138]{gimbaBoleVerbMorphology2000}} \\
\gll  a	ngòr-àakí-n-nì-n-gó \\
?	tie-\gl{vent}-\gl{vent}-3\gl{sg}.\gl{m}-\gl{vent}-\gl{pfv} \\
\glt `He will tie and bring (for) him.'

\ex\label{ex2:bole8}
\langinfo{Bole verb \textit{tí} `eat' with ventive suffix}{}{\cite[142]{gimbaBoleVerbMorphology2000}} \\
\gll tí-n-ní-n-tì \\
eat-\gl{vent}-3\gl{sg}.\gl{m}-\gl{vent}-\gl{tot} \\
\glt `He has eaten for him here.'
\z 

\subparagraph{Caused accompanied motion (CAM)}



The verb \textit{àlā} `carry, take from one place to another' is said to have a ventive form \textit{èlen} `bring' \citep{gimbaBoleenglishhausaDictionaryAlhaji2004}.

\subparagraph{Buy-sell verbs}



The verb \textit{gójju}  is translated `buy' in verbal form and `trading' in nominal form. The ventive form of the verb, \textit{gòjùttù}, is translated `sell' \citep[36]{gimbaBoleenglishhausaDictionaryAlhaji2004}.

\paragraph*{Galambu (gala1264)} 

The limited information available about Galambu grammar is from a chapter in a book of short grammar sketches \citep{schuhBoleTangaleLanguagesBauchi1978}. 

\subparagraph{Directional verbs}~



\begin{table}[H]
	\caption{Galambu directional verbs} 
	\label{tab:galaV}
	\begin{tabular}{lll} % add l for every additional column or remove as necessary
		\lsptoprule
		DIR & verb & gloss   \\ %table header
		\midrule
		ITIVE & \textit{nj-áalà} & `go' \\
		VENT & \textit{ndw-áalà} & `come' \\
		INTO&  \textit{ry-àalá} & `go in' \\
		OUT & \textit{páz-àalá} & `go out' \\
		\lspbottomrule
	\end{tabular}
\end{table}

\subparagraph{Directional extensions}

\citet[67]{schuhBoleTangaleLanguagesBauchi1978} reports that the Galambu speakers he worked with ``maintained that Galambu did not overtly express a ventive/neutral distinction.'' 

\subparagraph{Caused accompanied motion (CAM)}



Galambu has directed caused accompanied motion verbs derived from directional verbs ``by means of a suffix \textit{-z-} (< *\textit{t})'' \citep[76]{schuhBoleTangaleLanguagesBauchi1978}, as in (\ref{ex2:gala1}) and (\ref{ex2:gala2}). No other examples are given of this morpheme, but \citet[309]{schuhChadicCornucopia2017} proposes *\textit{t} as a proto-from of causative morphology in Bole-Tangale languages. 

\ea\label{ex2:gala1}
\langinfo{Gala itive caused accompanied motion}{}{\cite[76]{schuhBoleTangaleLanguagesBauchi1978}} \\
\gll  shì njíi-z-áalà \\
3\gl{sg} go-??-\gl{pfv} \\
\glt `He took.'

\ex\label{ex2:gala2}
\langinfo{Gala ventive caused accompanied motion}{}{\cite[76]{schuhBoleTangaleLanguagesBauchi1978}} \\
\gll  shì ndóo-z-áalà \\
3\gl{sg} come-??-\gl{pfv} \\
\glt `He brought.'
\z 
\subparagraph{Buy-sell verbs}



The verb \textit{bày-áalà} `buy' is similar in form to \textit{bàyáa kə̀rí} `sell' \citep[82]{schuhBoleTangaleLanguagesBauchi1978}, but it is not clear what the function of \textit{kə̀rí} is elsewhere in the language. 

\paragraph*{Gera (gera1246)} \label{gera}

The limited information available about the grammar of Gera is from a chapter in a book of short grammar sketches \citep{schuhBoleTangaleLanguagesBauchi1978}. 

\subparagraph{Directional verbs}~



\begin{table}[H]
	\caption{Gera directional verbs} 
	\label{tab:geraV}
	\begin{tabular}{lll} % add l for every additional column or remove as necessary
		\lsptoprule
		DIR & verb & gloss   \\ %table header
		\midrule
		ITIVE & \textit{ndíi-mí} & `go' \\
		VENT & \textit{ndûu-mí} & `come' \\
		INTO& \textit{ríi-mí} & `go in' \\
		OUT & \textit{fíɗ-mí} & `go out' \\
		\lspbottomrule
	\end{tabular}
\end{table}

\subparagraph{Directional extensions}~



\begin{table}[H]
	\caption{Gera directional morphology}
	\label{tab:gera}
	\begin{tabular}{llll} % add l for every additional column or remove as necessary
		\lsptoprule
		forms & source gloss & direction  & other functions  \\ %table header
		\midrule
		L tone  &  ventive  &   \textsc{vent} &   \textsc{subs.vent}  \\
		\lspbottomrule
	\end{tabular}
\end{table} 

\citet[99--100]{schuhBoleTangaleLanguagesBauchi1978} gives several examples of ventive verbs and comments on the morphological form: ``[Gera] does have a ventive, but the ventive and neutral forms are neutralized in many cases. For disyllabic verbs the ventive is shown by low tone on the verb. Therefore, the only cases where the ventive is distinct from the neutral form is for high tone CVC- verbs with singular subjects.''  These ventive forms can express ventive directional meaning when combined with a motion verb, as in (\ref{ex2:gera1}). 

\ea\label{ex2:gera1}
\langinfo{Gera verb \textit{rôomí} `enter' in ventive form}{}{\cite[100]{schuhBoleTangaleLanguagesBauchi1978}} \\
\gll sì rôomí \\
3\gl{sg} enter.\gl{vent}	\\
\glt `He came in.'
\z 

There are just two non-motion verbs in ventive form in the data. In one example, shown in (\ref{ex2:gera2}), the translation expresses subsequent ventive (caused accompanied) motion. 

\ea\label{ex2:gera2}
\langinfo{Gera verb \textit{zòomí} `buy' in ventive form}{}{\cite[100]{schuhBoleTangaleLanguagesBauchi1978}} \\
\gll sì	zòomí \\
3\gl{sg}	buy.\gl{vent}	 \\
\glt `He bought (and brought).'
\z 

The other non-motion verb shown in a ventive form is glossed ‘see’, shown in (\ref{ex2:gera3}). This form is said to be ambiguous between a ventive and non-ventive form. In this case, the free translation does not make explicit what the meaning of the ventive form is.

\ea\label{ex2:gera3}
\langinfo{Gera verb \textit{nêemí} `see' in ventive form}{}{\cite[100]{schuhBoleTangaleLanguagesBauchi1978}} \\
\gll sì	nêemí \\
3\gl{sg}	see(.\gl{vent}) \\
\glt `They saw.'
\z 

\subparagraph{Caused accompanied motion (CAM)}



It appears that a verb glossed `catch' can express ventive caused accompanied motion when in the ventive form, as in (\ref{ex2:gera4}). Another verb, glossed `take' and `carry', also expresses ventive CAM in the ventive form, as in (\ref{ex2:gera5}).

\ea\label{ex2:gera4}
\langinfo{Gera verb \textit{tàwmí} `catch' in ventive form}{}{\cite[100]{schuhBoleTangaleLanguagesBauchi1978}} \\
\gll sì	tàwmí \\
3\gl{sg}	catch.\gl{vent} \\
\glt `He caught (and brought it).'

\ex\label{ex2:gera5}
\langinfo{Gera verb \textit{gìɗmí} `take, carry' in ventive form}{}{\cite[100]{schuhBoleTangaleLanguagesBauchi1978}} \\
\gll  sì	gìɗmí \\
3\gl{sg}	carry.\gl{vent} \\
\glt `He carried (it) here.'
\z 

\subparagraph{Buy-sell verbs}



The verb \textit{zôo-mí} can mean `buy' or `sell' \citep[121]{schuhBoleTangaleLanguagesBauchi1978}. In the few examples where this verb appears in a ventive form it is translated `buy', as in (\ref{ex2:gera2}), but the ventive appears to express subsequent motion, rather than serving to distinguish the buying/selling interpretations of the verb.

\paragraph*{Geruma (geru1240)} 

The limited information available about the grammar of Geruma is from a chapter in a book of short grammar sketches \citep{schuhBoleTangaleLanguagesBauchi1978}. 

\subparagraph{Directional verbs}~



\begin{table}[H]
	\caption{Geruma directional verbs} 
	\label{tab:geruV}
	\begin{tabular}{lll} % add l for every additional column or remove as necessary
		\lsptoprule
		DIR & verb & gloss   \\ %table header
		\midrule
		VENT & \textit{ni-ŋŋ-álà} & `come' \\
		ITIVE & \textit{yóo-là} & `go' \\
		INTO& \textit{ríi-lá} & `go in' \\
		OUT & \textit{fát-álà} & `go out' \\
		\lspbottomrule
	\end{tabular}
\end{table}

The \textit{ŋŋ} morpheme in the verb `come' is not explained, and this morpheme does not appear in all the conjugations of the verb. 

\subparagraph{Directional extensions}~



\begin{table}[H]
	\caption{Geruma directional morphology}
	\label{tab:geru}
	\begin{tabular}{llll} % add l for every additional column or remove as necessary
		\lsptoprule
		forms & source gloss & direction  & other functions  \\ %table header
		\midrule
		\textit{-ìŋ}  & ventive   &   \textsc{vent} &     \\
		\lspbottomrule
	\end{tabular}
\end{table} 

The only relevant information available on the ventive suffix in Geruma is a brief description of its form: ``When no object is expressed, the ventive is formed by (a) giving low tone to the verb root, (b) adding a suffix -ìŋ'' \citep[127]{schuhBoleTangaleLanguagesBauchi1978}. This is exemplified with one verb root, e.g., \textit{nà gyátálà} ‘I carried (it)’ compared with \textit{nà gyàtìŋálà} ‘I brought (it)’. There is discussion of the form of the ventive when co-occuring with pronominals, but no further discussion of its functions.

\subparagraph{Caused accompanied motion (CAM)}



As seen above, the ventive form of the verb `carry' can express directed caused accompanied motion. No additional data on directed CAM in Geruma are available. 

\subparagraph{Buy-sell verbs}



The verb \textit{jáw-álà} can mean `buy' or `sell' \citep[136]{schuhBoleTangaleLanguagesBauchi1978}. There is no data available concerning how the interpretations might be distinguished. 

\paragraph*{Giiwo (giiw1236)} \label{giiw} 

The limited information available about the grammar of Giiwo is from a chapter in a book of short grammar sketches \citep{schuhBoleTangaleLanguagesBauchi1978}. 

\subparagraph{Directional verbs}~



\begin{table}[H]
	\caption{Giiwo directional verbs} 
	\label{tab:giiwV}
	\begin{tabular}{lll} % add l for every additional column or remove as necessary
		\lsptoprule
		DIR & verb & gloss   \\ %table header
		\midrule
		ITIVE & \textit{yów-wò} & `go' \\
		VENT & \textit{nón-kò} & `come' \\
		INTO & \textit{ríi-wò}  & `go in' \\
		OUT &  \textit{fák-kò}  & `go out' \\
		\lspbottomrule
	\end{tabular}
\end{table}

\subparagraph{Directional extensions}~



\begin{table}[H]
	\caption{Giiwo directional morphology}
	\label{tab:giiw}
	\begin{tabular}{llll} % add l for every additional column or remove as necessary
		\lsptoprule
		forms & source gloss & direction  & other functions  \\ %table header
		\midrule
		\textit{-(í)n}  &  ventive &   \textsc{vent}  &   \textsc{subs.vent}  \\
		\lspbottomrule
	\end{tabular}
\end{table} 

\citet[39]{schuhBoleTangaleLanguagesBauchi1978} includes a short description of the ventive form in Giiwo: ``The ventive extension is identical to the plural subject suffix in form, viz. \textit{-n} on monosyllabic verbs, \textit{-in} on all others. All verbs except monosyllabic take low tone. Since the plural subject suffix and the ventive extension are mutually exclusive, there can be no neutral/ventive distinction when the subject is plural.'' Six examples of ventive forms are given. These include intransitive and transitive motion verbs with ventive directional interpretation, as in (\ref{ex2:giiw1}) and (\ref{ex2:giiw2}). 

\ea\label{ex2:giiw1}
\langinfo{Giiwo verb \textit{rí} `enter' with ventive suffix}{}{\cite[39]{schuhBoleTangaleLanguagesBauchi1978}} \\
\gll shì	rí-n-kò \\
3\gl{sg}	enter-\gl{vent}-\gl{pfv} \\
\glt `He came in.'

\ex\label{ex2:giiw2}
\langinfo{Giiwo verb \textit{bàr} `give' with ventive suffix}{}{\cite[39]{schuhBoleTangaleLanguagesBauchi1978}} \\
\gll  shì	bàr-ín-kò \\
3\gl{sg}	give-\gl{vent}-\gl{pfv} \\
\glt `He gave (it) here.'
\z 

\newpage
Note that in (\ref{ex2:giiw3}) it is not clear that any particular argument of the verb is the figure on the path of motion, but rather the ventive meaning refers to the orientation of the stabbing event (or the path of motion of the unexpressed instrument). 

\ea\label{ex2:giiw3}
\langinfo{Giiwo verb \textit{tùɗ} `stab' with ventive suffix}{}{\cite[39]{schuhBoleTangaleLanguagesBauchi1978}} \\
\gll  shì	tùɗ-ín-kò \\
3\gl{sg}	stab-\gl{vent}-\gl{pfv} \\
\glt `He stabbed (in this direction).'
\z 

One example indicates that it is also possible for the ventive suffix to be used to express subsequent (caused accompanied) motion, shown in (\ref{ex2:giiw4}).

\ea\label{ex2:giiw4}
\langinfo{Giiwo verb \textit{màt} `shoot' with ventive suffix}{}{\cite[39]{schuhBoleTangaleLanguagesBauchi1978}} \\
\gll  mù	màt-ín-kò \\
1\gl{pl}	shoot-\gl{vent}-\gl{pfv} \\
\glt `We shot (and brought).' or `We shot.'
\z 

\subparagraph{Caused accompanied motion (CAM)}



The verb \textit{ɗé} `get' can have a ventive caused accompanied motion meaning in its ventive form, as in (\ref{ex2:giiw5}). 

\ea\label{ex2:giiw5}
\langinfo{Giiwo verb \textit{ɗé} `get' with ventive suffix}{}{\cite[39]{schuhBoleTangaleLanguagesBauchi1978}} \\
\gll  shì	ɗé-n-kò \\
3\gl{sg}	get-\gl{vent}-\gl{pfv} \\
\glt `He got (it and brought it).'
\z 

\subparagraph{Buy-sell verbs}



The verb \textit{bàayú-wò} ‘buy’ is similar in form to \textit{bàyíté} ‘sell’. \citet[52]{schuhBoleTangaleLanguagesBauchi1978} includes a note ``(? form)'' next to the verb \textit{bàyíté} ‘sell’ in the Giiwo wordlist, indicating that he was not able to identify the apparent suffix that distinguishes the two verbs. 

\paragraph*{Ngamo (ngam1282)} \label{ngam}

Information on the grammar of Ngamao is from several sources. Schuh published a journal article on the morphology \citep{schuhDegeminationCompensatoryLengthening2005} and included Ngamo in sections of his book on Chadic languages \citep[332--353]{schuhChadicCornucopia2017} in addition to making unpublished field notes and drafts available online \citep{schuhNgamoVerbalSystem2004,schuhGudiNgamoVerb2010}. There is also a chapter on Ngamo verbs in an edited volume \citep{ibriszimowVerbNgamo2006}.

\subparagraph{Directional verbs}~



\begin{table}[H]
	\caption{Ngamo directional verbs} 
	\label{tab:ngamV}
	\begin{tabular}{lll} % add l for every additional column or remove as necessary
		\lsptoprule
		DIR & verb & gloss   \\ %table header
		\midrule
		MOTION & \textit{ndu} & `go/come' \\
		UP & \textit{ɗâ} & `climb, mount' \\
		DOWN & \textit{ùr-kô} & `get down, descend' \\
		INTO& \textit{rukò} & `enter' \\
		OUT & \textit{hàtâ/fàtâ} & `go out' \\
		\lspbottomrule
	\end{tabular}
\end{table}

\citet[336]{schuhChadicCornucopia2017} glosses the verb root \textit{ndu-} as `go/come'. The verb root can also have a ventive form in which case it can only express ventive directional motion. \citet{schuhNgamoVerbalSystem2004} notes that the ventive paradigm of this verb includes several irregular forms.

\subparagraph{Directional extensions}~



\begin{table}[H]
	\caption{Ngamo directional morphology}
	\label{tab:ngam}
	\begin{tabular}{llll} % add l for every additional column or remove as necessary
		\lsptoprule
		forms & source gloss & direction  & other functions  \\ %table header
		\midrule
		\textit{-no, -n, -t}  &  ventive  &  \textsc{vent} &   \textsc{subs.vent}, \textsc{ben}  \\
		\lspbottomrule
	\end{tabular}
\end{table} 

\citet{schuhDegeminationCompensatoryLengthening2005} and \citet[332--353]{schuhChadicCornucopia2017} describe the morphology of Ngamo verbs including a few translations of ventive forms. The function of the ventive is defined as ``indicating action initiated at a distance with effect at the place of reference.'' Schuh’s \citeyearpar{schuhNgamoVerbalSystem2004,schuhGudiNgamoVerb2010} field notes and \citet{ibriszimowVerbNgamo2006} also include ventive forms, but without translations. In (\ref{ex2:ngam1}) the ventive form of an intransitive motion verb has a directional interpretation. 

\ea\label{ex2:ngam1}
\langinfo{Ngamo verb \textit{hete} `exit' with ventive suffix}{}{\cite[335]{schuhChadicCornucopia2017}} \\
\gll baɗi wuya, sai shapsu hete-tu a ramte zugonsu.	 \\
?? ?? ?? ?? exit.\gl{sbjv}-\gl{vent} ?? ?? ??		\\
\glt `When night ended, both of them came out and readied themselves.'
\z 

%In (\ref{ex2:ngam2}), the ventive form of the transitive verb `call' has a directional interpretation (with the patient as the figure on a path of motion).
%\ea\label{ex2:ngam2}
%\langinfo{Ngamo verb \textit{èeʃè} `call' with ventive suffix}{}{\cite[3]{schuhDegeminationCompensatoryLengthening2005}} \\
%\gll èeʃè-n-tò \\
%call-\gl{vent}-3\gl{sg}.\gl{f}.\gl{dat} \\
%\glt `He called hither for her.'
%\z 

\newpage
In (\ref{ex2:ngam3}), the ventive form of the verb `tie' has a subsequent caused accompanied motion interpretation. However, when an indirect object pronoun occurs with the ventive suffix, as in (\ref{ex2:ngam6}), the ventive form has a benefactive interpretation.

\ea\label{ex2:ngam3}
\langinfo{Ngamo verb \textit{ŋgàr} `tie' with ventive suffix}{}{\cite[3]{schuhDegeminationCompensatoryLengthening2005}} \\
\gll ŋgàr-nô \\
tie-\gl{vent} \\
\glt `He tied and brought.'

\ex\label{ex2:ngam6}
\langinfo{Ngamo verb \textit{ŋgàr} `tie' with ventive suffix}{}{\cite[3]{schuhDegeminationCompensatoryLengthening2005}} \\
\gll à ŋgàr-ii-tò \\
?? tie-\gl{vent}-3\gl{sg}.\gl{f}.\gl{dat}	\\
\glt `That he tie for her.'
\z 

\subparagraph{Caused accompanied motion (CAM)}



Very limited information is available about directed CAM in Ngamo. The causative forms of the upward and downward directional verbs can have a directed CAM interpretation, as in (\ref{ex2:ngam4}) and (\ref{ex2:ngam5}), but no examples were found of causative forms of the basic motion verb \textit{ndu-} `go/come'. 

\ea\label{ex2:ngam4}
\langinfo{Ngamo verb \textit{ɗâ} `climb, mount' with causative suffix}{}{\cite[333]{schuhChadicCornucopia2017}} \\
\gll ɗā̀-t-â  \\
climb-\gl{caus}-?? \\
\glt `put onto'

\ex\label{ex2:ngam5}
\langinfo{Ngamo verb \textit{ùr} `get down, descend' with causative suffix}{}{\cite[333]{schuhChadicCornucopia2017}} \\
\gll ùr-t-â \\
descend-\gl{caus}-?? \\
\glt `take down a load'
\z

\subparagraph{Buy-sell verbs}


\sloppy
\citet[332]{schuhChadicCornucopia2017} describes the verb \textit{bò’ota} `sold' as the caus\-a\-tive form of the verb \textit{kàja} `bought', but causative forms are normally marked by the affix \textit{-t}. This appears to be an over-analysis based on the use of causative forms in other languages. The two verbs are here treated as separate, unrelated forms.
\fussy

\paragraph*{Bure (bure1242)}

Information on the grammar of Bure is from a grammar sketch \citep{baticGrammaticalSketchBure2014}. 

\subparagraph{Directional verbs}~



\begin{table}[H]
	\caption{Bure directional verbs} 
	\label{tab:bureV}
	\begin{tabular}{lll} % add l for every additional column or remove as necessary
		\lsptoprule
		DIR & verb & gloss   \\ %table header
		\midrule
		ITIVE & \textit{yo'} & `go, leave' \\
		INTO & \textit{ri'-} & `enter' \\
		OUT & \textit{pat-} & `go out' \\
		DOWN & \textit{ndur-} & `descend' \\
		\lspbottomrule
	\end{tabular}
\end{table}

Where English `come' appears in translations, the corresponding Bure verb is \textit{maɗ-} `return' with a ventive extension. 

\subparagraph{Directional extensions}~



\begin{table}[H]
	\caption{Bure directional morphology}
	\label{tab:bure}
	\begin{tabular}{llll} % add l for every additional column or remove as necessary
		\lsptoprule
		forms & source gloss & direction  & other functions  \\ %table header
		\midrule
		\textit{-ín}  &  \textsc{vt}  &  \textsc{vent} &  \textsc{subs.vent}  \\
		\lspbottomrule
	\end{tabular}
\end{table} 

\citet[70]{baticGrammaticalSketchBure2014} briefly describes a ventive extension \textit{-ín}. Not many examples are given, but the suffix is shown to have ventive directional meaning with some intransitive motion verbs, as in (\ref{ex2:bure1}) and (\ref{ex2:bure2}). 

\ea\label{ex2:bure1}
\langinfo{Bure verb \textit{rí'} `enter' with ventive suffix}{}{\cite[64]{baticGrammaticalSketchBure2014}} \\
\gll kà máa bán-rù mìmmíŋ sú-ntì rí'-ín-kò	\\
\gl{prf}:2\gl{sg}.\gl{m} \gl{neg} know-\gl{neg} men 3\gl{pl}-\gl{rel} enter-\gl{vent}-\gl{prf}		\\
\glt `You don’t know the men who came in.'

\ex\label{ex2:bure2}
\langinfo{Bure verb \textit{màɗ} `return' with ventive suffix}{}{\cite[70]{baticGrammaticalSketchBure2014}} \\
\gll màɗ-ín-kò	\\
return-\gl{vent}-\gl{prf}		\\
\glt `We came back.'
\z 

With one non-motion verb there is an apparent subsequent (caused accompanied) motion translation, shown in (\ref{ex2:bure3}).

\ea\label{ex2:bure3}
\langinfo{Bure verb \textit{mòr} `steal' with ventive suffix}{}{\cite[70]{baticGrammaticalSketchBure2014}} \\
\gll mòr-ín-kò	\\
steal-\gl{vent}-\gl{prf}	\\
\glt `We stole and brought [it].'
\z 

\subparagraph{Caused accompanied motion (CAM)}



The ventive suffix \textit{-ín} is shown with several verbs expressing caused accompanied motion. The verb stems are often glossed `take', but other examples of these same verb roots show that they are able to express CAM without the ventive suffix, and so are inherently motion verbs. 

\subparagraph{Buy-sell verbs}



\citet[99]{baticGrammaticalSketchBure2014} includes one example with the verb \textit{jò} `buy' but the expression for the notion `sell' is not given. 

\paragraph*{Nyam (nyam1285)} \label{nyam}

Information on the grammar of Nyam is from a doctoral dissertation \citep{andreasGrammatischeBeschreibungNyam2012}. 

\subparagraph{Directional verbs}~



\begin{table}[H]
	\caption{Nyam directional verbs} 
	\label{tab:nyamV}
	\begin{tabular}{lll} % add l for every additional column or remove as necessary
		\lsptoprule
		DIR & verb & gloss   \\ %table header
		\midrule
		ITIVE & \textit{sí} & `\textit{gehen} [go]' \\
		VENT & \textit{tò} & `\textit{kommen} [come]' \\
		UP & \textit{ʔɔ̀nd} & `\textit{aufsteigen} [ascend]' \\
		\lspbottomrule
	\end{tabular}
\end{table}

\subparagraph{Directional extensions}~



\citet[212]{andreasGrammatischeBeschreibungNyam2012} notes that ``\textit{Einige traditionelle Merkmale von Bole-Tangale-Sprachen wie beispielsweise ein binäres Aspektsystem bei den Verben, logophorische Pronomina oder auch} ,Intransitive Copy Pronouns‘ (ICPs) \textit{sowie Verbalerweiterungen wie z.B. Ventiv oder Destinativ besitzt das Nyam nicht}. [Nyam does not have some traditional features of Bole-Tangale languages such as a binary aspect system for verbs, logophoric pronouns or `Intransitive Copy Pronouns' (ICPs) and verbal extensions such as ventive or destinative.]''

\subparagraph{Caused accompanied motion (CAM)}



Directed caused accompanied motion may be expressed by using a motion verb \textit{tò} `\textit{kommen} [come]' with a comitative prepositional phrase, as can be seen by comparing (\ref{ex2:nyam1}) and (\ref{ex2:nyam2}).

\ea\label{ex2:nyam1}
\langinfo{Nyam verb \textit{tò} `come'}{}{\cite[127]{andreasGrammatischeBeschreibungNyam2012}} \\
\gll mùdùk tò-wá \\
woman come-\gl{prf} \\
\glt `The woman has come.' [German: `\textit{Die Frau ist gekommen}.']

\ex\label{ex2:nyam2}
\langinfo{Nyam verb \textit{tò} `come' with comitative prepositional phrase}{}{\cite[127]{andreasGrammatischeBeschreibungNyam2012}} \\
\gll mùdùk tò-wá yà kɔ̀lɔ́ŋ \\
woman come-\gl{prf} with food \\
\glt `The woman has brought food.' [German: `\textit{Die Frau hat das Essen gebracht}.']
\z 

\subparagraph{Buy-sell verbs}



In addition to the use of comitative phrases to express CAM, as mentioned above, there are a number of verbs in which a comitative prepositional phrase gives an interpretation that \citet[127]{andreasGrammatischeBeschreibungNyam2012} describes as causative. This includes \textit{ʔilg-} `awaken' (becomes `wake (someone)'), \textit{so} `eat' (becomes `feed') and \textit{kemd-} `buy' (becomes `sell'), as in (\ref{ex2:nyam3}). 

\ea\label{ex2:nyam3}
\langinfo{Nyam verb \textit{kemd} `buy' with comitative prepositional phrase}{}{\cite[127]{andreasGrammatischeBeschreibungNyam2012}} \\
\gll mùdùk kémd-ì yà brédì \\
woman sell-\gl{hab} with bread \\
\glt `The woman usually sells bread.' [German: `\textit{Die Frau verkauft gewöhnlich Brot}.']
\z 

\subsubsection{Tangalic, West Chadic A2 (4/7)}

Information on the grammar of Tangalic languages is available for Kanakuru, Kushi, Pero and Tangale. \citet{legerGrammatikKwamiSpracheNordostnigeria1994} describes the grammar of Kwaami (kwaa1269) including a suffix labeled ``\textit{Ventiv}''. However, the examples given do not show any clear cases of directional meaning, and without glosses or a lexicon, it is not possible to gain any insight into the lexical semantics of the language either. No information was found for Piya-Kwonci (piya1245) or Kutto (kutt1236).

\paragraph*{Kanakuru (dera1248)} \label{kana}

Kanakuru (or Dera) is described in a short book that focuses on a transformational analysis of the grammar \citep{newmanKanakuruLanguage1974}. 

\subparagraph{Directional verbs}~



\begin{table}[H]
	\caption{Kanakuru directional verbs} 
	\label{tab:kanaV}
	\begin{tabular}{lll} % add l for every additional column or remove as necessary
		\lsptoprule
		DIR & verb & gloss   \\ %table header
		\midrule
		ITIVE & \textit{tál} & `go' \\
		VENT & \textit{dol} & `come' \\
		UP & \textit{yíll} & `rise, raise' \\
		DOWN &  \textit{kèeɓé/w-} & `get down' \\
		INTO& \textit{gàl} & `enter' \\
		OUT & \textit{púl} & `go out' \\
		\lspbottomrule
	\end{tabular}
\end{table}

\citet[131]{newmanKanakuruLanguage1974} notes that it is more common to express outward motion using the intransitive form of the labile verb \textit{pòrí} `take out, go out'.

\subparagraph{Directional extensions}~



\begin{table}[H]
	\caption{Kanakuru directional morphology}
	\label{tab:kana}
	\begin{tabular}{llll} % add l for every additional column or remove as necessary
		\lsptoprule
		forms & source gloss & direction  & other functions  \\ %table header
		\midrule
		\textit{-tə(ru)}  &  ventive  &   \textsc{vent} &   \textsc{subs.vent}, \textsc{prior.itive}, \textsc{ben}, \textsc{loc} \\
		\textit{bo}  &   ventive aux  &   \textsc{vent} &   \textsc{prior.itive}, \textsc{loc}  \\
		\lspbottomrule
	\end{tabular}
\end{table} 

The ventive suffix in Kanakuru is said to have three allomorphs depending on the phonological and morphological context. With intransitive and transitive motion verbs, the ventive suffix has a ventive directional interpretation, as in (\ref{ex2:kana1}) and (\ref{ex2:kana2}). 

\ea\label{ex2:kana1}
\langinfo{Kanakuru verb \textit{ma} `return' with ventive suffix}{}{\cite[73]{newmanKanakuruLanguage1974}} \\
\gll à	ma-tə-ni \\
\gl{pfv}	return-\gl{vent}-\gl{pfv} \\
\glt `He returned here.'

\ex\label{ex2:kana2}
\langinfo{Kanakuru verb \textit{àm} `pull' with ventive suffix}{}{\cite[74]{newmanKanakuruLanguage1974}} \\
\gll nà	àm-tə̀ wú \\
1\gl{sg}	pull-\gl{vent}	3\gl{pl}.\gl{obj} \\
\glt `I pulled them (here).'
\z 

With non-motion verbs, there are a few different interpretations given of the ventive suffix. Newman gives several examples of the ventive suffix expressing prior itive associated motion with a non-motion verb, as in (\ref{ex2:kana3}), (\ref{ex2:kana4}), (\ref{ex2:kana5}). Newman does not discuss the semantics of these examples, but the translations were confirmed by a Kanakuru speaker via social media.\footnote{Personal communication with Joel Solomon Douro Talum in the Facebook group National Association of Dera (Kanakuru) Student Union (NADSU).}

\ea\label{ex2:kana3}
\langinfo{Kanakuru verb \textit{ɗop} `tie' with ventive suffix}{}{\cite[7]{newmanKanakuruLanguage1974}} \\
\gll a	ɗop-təru \\
\gl{pfv}	tie-\gl{vent} \\
\glt `He went and tied it.'

\ex\label{ex2:kana4}
\langinfo{Kanakuru verb \textit{shak} `found' with ventive suffix}{}{\cite[7]{newmanKanakuruLanguage1974}} \\
\gll  a	shak-təru \\
\gl{pfv}	found-\gl{vent} \\
\glt `He went and founded it.'

\ex\label{ex2:kana5}
\langinfo{Kanakuru verb \textit{lok} `hang' with ventive suffix}{}{\cite[9]{newmanKanakuruLanguage1974}} \\
\gll  a	lok-təru \\
\gl{pfv}	hang-\gl{vent} \\
\glt `He (went) and hung it.'
\z 

A different meaning is found when the ventive form of a non-motion verb is followed by a dative pronominal suffix. In (\ref{ex2:kana6}) the ventive has benefactive meaning in regard to the dative pronominal. In contrast, the dative pronominal in (\ref{ex2:kana7}) (without a ventive suffix) has a malefactive or source interpretation.

\ea\label{ex2:kana6}
\langinfo{Kanakuru verb \textit{shit} `steal' with ventive suffix and pronominal suffix}{}{\cite[73]{newmanKanakuruLanguage1974}} \\
\gll à	shit-tə-no	dok \\
\gl{pfv}	steal-\gl{vent}-1\gl{sg}.\gl{dat}	horse \\
\glt `He stole a horse for me.'

\ex\label{ex2:kana7}
\langinfo{Kanakuru verb \textit{shit} `steal' with pronominal suffix}{}{\cite[73]{newmanKanakuruLanguage1974}} \\
\gll  à	shin-no	dok \\
\gl{pfv}	steal-1\gl{sg}.\gl{dat}	horse \\
\glt `He stole a horse from me.'
\z 

The translation in (\ref{ex2:kana8}) suggests that the ventive marker has a possible function of expressing the distance of a direct object from the deictic center. This function is not discussed by Newman.

\ea\label{ex2:kana8}
\langinfo{Kanakuru verb \textit{àl} `see' with ventive suffix}{}{\cite[74]{newmanKanakuruLanguage1974}} \\
\gll mə̀	àl-də̀	ré \\
1\gl{pl}	see-\gl{vent}	3\gl{sg}.\gl{f}.\gl{obj} \\
\glt `We saw her (at a distance).'
\z 

\citet[76]{newmanKanakuruLanguage1974} also presents a ``ventive equivalent'' auxiliary \textit{bo} which appears before a deverbal predicate in certain tenses where the ventive suffix is not available, as in (\ref{ex2:kana9}). Newman notes that with motion verbs, the use of \textit{bo} is parallel to the ventive forms in Kanakuru and Hausa. 

\ea\label{ex2:kana9}
\langinfo{Kanakuru verb \textit{ma} `return' with \textit{bo} auxiliary}{}{\cite[76]{newmanKanakuruLanguage1974}} \\
\gll shìi	bo	ma-ma	\\
3\gl{sg}	\gl{vent}.\gl{aux}	return-return	 \\
\glt `He is returning (coming back).'
\z 

Newman remarks that with non-motion verbs “a \textit{bo} construction was considered somewhat different from a corresponding ventive stem, the denotation being primarily action at some distance from the speaker rather than in the direction of or for the benefit of the speaker” \citep[76]{newmanKanakuruLanguage1974}. In these cases, the translations given are either expressing prior associated motion, as in (\ref{ex2:kana10}), or translocative `there', as in (\ref{ex2:kana11}). It is not clear why the translations vary in this regard, though it is noteworthy that the prior motion interpretations match the translations of the ventive suffix with non-motion verbs described above (although Newman does not comment on this parallel). 

\ea\label{ex2:kana10}
\langinfo{Kanakuru verb \textit{wupe} `sell' with \textit{bo} auxiliary}{}{\cite[76]{newmanKanakuruLanguage1974}} \\
\gll àm	bo	wupe	kurei \\
1\gl{pl}	\gl{vent}.\gl{aux}	sell	corn \\
\glt `We will (go) sell the corn.'

\ex\label{ex2:kana11}
\langinfo{Kanakuru verb \textit{wuɗ} `herd' with \textit{bo} auxiliary}{}{\cite[76]{newmanKanakuruLanguage1974}} \\
\gll jítò	bo	wuɗ-mai \\
3\gl{sg}.\gl{f}.\gl{pst}.\gl{cont}	\gl{vent}.\gl{aux}	herd-3\gl{sg}.\gl{m}.\gl{obj} \\
\glt `She used to herd it (there).'
\z 

\subparagraph{Caused accompanied motion (CAM)}



The caused accompanied motion (CAM) verb \textit{gət} `carry' can take a ventive directional suffix, as in (\ref{ex2:kana12}).  The ventive form of the verb \textit{ko} ‘catch’ also has a ventive CAM interpretation, as in (\ref{ex2:kana13}).

\ea\label{ex2:kana12}
\langinfo{Kanakuru verb \textit{gət} `carry' with ventive suffix}{}{\cite[7]{newmanKanakuruLanguage1974}} \\
\gll a	gət-təru \\
\gl{pfv}	carry-\gl{vent} \\
\glt `He carried it here.'

\ex\label{ex2:kana13}
\langinfo{Kanakuru verb \textit{ko} `catch' with ventive suffix}{}{\cite[73]{newmanKanakuruLanguage1974}} \\
\gll à	ko-ru \\
\gl{pfv}	catch-\gl{vent} \\
\glt `He caught it (there and brought it here).'
\z

There are also labile motion verbs that can express directed CAM in their transitive form: \textit{pòrí} `take out, go out' and \textit{kèeɓé}/\textit{-w-} `get down, take down'. Other motion verbs can take the associative suffix \textit{-nú} which transitivizes them and gives them a CAM interpretation, as in (\ref{ex2:kana14}).

\ea\label{ex2:kana14}
\langinfo{Kanakuru verb \textit{dol} `come' with associative suffix}{}{\cite[25]{newmanKanakuruLanguage1974}} \\
\gll  à do-to-nu \\
\gl{pfv} come-3\gl{sg}.\gl{f}-\gl{asoc} \\
\glt `She brought (it).'
\z 

\subparagraph{Buy-sell verbs}



There are two separate verbs in Kanakuru, \textit{ɗìbəré} `buy' and \textit{wùpé} `sell' \citep[124, 134]{newmanKanakuruLanguage1974}.

\paragraph*{Kushi (kush1236)} \label{kush}

Information on the grammar of Kushi is from a grammar sketch of the verbal system included as a chapter in an edited volume \citep{baticVerbClassesTAM2019}.

\subparagraph{Directional verbs}~

\begin{table}[H]
	\caption{Kushi directional verbs} 
	\label{tab:kushV}
	\begin{tabular}{lll} % add l for every additional column or remove as necessary
		\lsptoprule
		DIR & verb & gloss   \\ %table header
		\midrule
		ITIVE & \textit{kʰɔ̀ɔ} & `go' \\
		VENT & \textit{wàr-} & `come' \\
		INTO& \textit{rè} & `enter' \\
		OUT & \textit{pɛ̀rɔ̀} & `go out' \\
		\lspbottomrule
	\end{tabular}
\end{table}

\newpage
\subparagraph{Directional extensions}~

\begin{table}[H]
	\caption{Kushi directional morphology}
	\label{tab:kush}
	\begin{tabular}{llll} % add l for every additional column or remove as necessary
		\lsptoprule
		forms & source gloss & direction  & other functions \\ %table header
		\midrule
		\textit{-na, -ru}  &  \textsc{vent}  &  \textsc{vent} &     \\
		\lspbottomrule
	\end{tabular}
\end{table} 

The ventive suffix in Kushi has various forms depending on the TAM of the verb. \citet{baticVerbClassesTAM2019} present examples of the suffix with intransitive motion verbs where it has a directional function, as in (\ref{ex2:kush1}) and (\ref{ex2:kush2}).

\ea\label{ex2:kush1}
\langinfo{Kushi verb \textit{ʔàm} `climb' with ventive suffix}{}{\cite[78]{baticVerbClassesTAM2019}} \\
\gll shɪ̀ɪ ʔàm-nà shɪ̀k-nɔ̀	\\
3\gl{sg}.\gl{m} climb-\gl{pfv}.\gl{vent} body-1\gl{sg}.\gl{poss}		\\
\glt `He climbed towards me.'

\ex\label{ex2:kush2}
\langinfo{Kushi verb \textit{pì} `exit' with ventive suffix}{}{\cite[79]{baticVerbClassesTAM2019}} \\
\gll pì-rù	\\
exit-\gl{vent}		\\
\glt `Come out.'
\z 

\subparagraph{Caused accompanied motion (CAM)}



No information available. 

\subparagraph{Buy-sell verbs}



One example in \citet[75]{baticVerbClassesTAM2019} includes the verb \textit{pɪ̀l-ànì} `he is buying' but the expression for the notion `sell' in Kushi was not found. 

\paragraph*{Pero (pero1241)} \label{pero}

Information on the grammar of Pero is from a reference grammar \citep{frajzyngierGrammarPero1989}. 

\largerpage
\subparagraph{Directional verbs}~



\begin{table}[H]
	\caption{Pero directional verbs} 
	\label{tab:peroV}
	\begin{tabular}{lll} % add l for every additional column or remove as necessary
		\lsptoprule
		DIR & verb & gloss   \\ %table header
		\midrule
		MOTION & \textit{ẃaatò} & `go, come' \\
		UP & \textit{ámbó} & `go up, come up' \\
		DOWN & \textit{cáa} & `descend' \\
		INTO& \textit{rí} & `enter' \\
		OUT & \textit{pétò} & `go out, leave'  \\
		\lspbottomrule
	\end{tabular}
\end{table}

\subparagraph{Directional extensions}~



\begin{table}[H]
	\caption{Pero directional morphology}
	\label{tab:pero}
	\begin{tabular}{llll} % add l for every additional column or remove as necessary
		\lsptoprule
		forms & source gloss & direction  & other functions \\ %table header
		\midrule
		\textit{-na, -tu}  &  \textsc{vent}  &  \textsc{vent} &  \textsc{subs.vent}  \\
		\lspbottomrule
	\end{tabular}
\end{table} 

Pero is described by \citet[95]{frajzyngierGrammarPero1989} as having two ventive verbal suffixes which vary according to aspect: the completive ventive form \textit{-na} and the non-completive ventive form \textit{-tu}. The general motion verb \textit{ẃaatò} has an itive `go' default interpretation, and with a ventive suffix it means `come', as in (\ref{ex2:pero1}). 

\ea\label{ex2:pero1}
\langinfo{Pero verb \textit{ẃaatò} `go/come' with ventive suffix}{}{\cite[88]{frajzyngierGrammarPero1989}} \\
\gll  wát-tù	mìnáa-nò	 \\
go/come-\gl{vent}	house-1\gl{sg}	 \\
\glt `Come to my house.'
\z

Ventive directional meaning can also be seen in the contrasting pair in (\ref{ex2:pero2}) and (\ref{ex2:pero3}). In (\ref{ex2:pero2}), the motion ‘climb’ is understood to be on a path toward the deictic center (the location of the speaker). In contrast, in (\ref{ex2:pero3}), the same verb without a ventive suffix is translated as motion away from the deictic center. 

\ea\label{ex2:pero2}
\langinfo{Pero verb \textit{ámb} `climb' with ventive suffix}{}{\cite[95]{frajzyngierGrammarPero1989}} \\
\gll  ámb-ínà	tì	pókáyà \\
climb-\gl{compl}.\gl{vent}	\gl{prep}	west \\
\glt `He came from the west. (The speaker is at the place to which x arrived.)'

\ex\label{ex2:pero3}
\langinfo{Pero verb \textit{ámb} `climb' without ventive suffix}{}{\cite[95]{frajzyngierGrammarPero1989}} \\
\gll  ámbì-kò	itì	pókáyà \\
climb-\gl{compl}	\gl{prep}	west \\
\glt `He came from the west. (The speaker is not at the place to which x arrived.)'
\z 

\citet[94]{frajzyngierGrammarPero1989} describes subsequent associated motion as another possible function of the ventive suffix: ``There are two crucial components of the function of the ventive. One is that the action, process or event takes place outside of the place of speech or some other previously defined place, and the other is that there be a subsequent movement toward the place of speech.'' Only one unambiguous example of subsequent associated motion is given, shown in (\ref{ex2:pero4}). However, there are two other examples of the same verb in a ventive form, as in (\ref{ex2:pero5}), where the given translation does not indicate any motion meaning related to the ventive suffix \citep[see also][242]{frajzyngierGrammarPero1989}.

\ea\label{ex2:pero4}
\langinfo{Pero verb \textit{cúg} `fall' with ventive suffix}{}{\cite[94]{frajzyngierGrammarPero1989}} \\
\gll  cúg-ínà	tù	púccù \\
fall-\gl{compl}.\gl{vent}	\gl{prep}	there	 \\
\glt `He fell there and subsequently came here.'

\ex\label{ex2:pero5}
\langinfo{Pero verb \textit{cúg} `fall' with ventive suffix}{}{\cite[241]{frajzyngierGrammarPero1989}} \\
\gll nì-íkkà	wáatò	mínà	nì-n-cúg-ínà \\
1\gl{sg}-\gl{prog}	go/come	home	1\gl{sg}-and-fall-\gl{compl}.\gl{vent} \\
\glt `I was going home and I fell down.'
\z

There are at least 20 other verbs in a ventive form for which no motion-related meaning in given in the translation, as in (\ref{ex2:pero6}) and (\ref{ex2:pero7}).

\ea\label{ex2:pero6}
\langinfo{Pero verb \textit{tákl} `rub' with ventive suffix}{}{\cite[133]{frajzyngierGrammarPero1989}} \\
\gll  n-tákl-ínà	kòngò-ì	n-ɓál-éetò \\
and-rub-\gl{vent}	stomach-\gl{def}	and-burst-3\gl{f} \\
\glt `And the stomach was rubbed and it burst.'

\ex\label{ex2:pero7}
\langinfo{Pero verb \textit{áɗɗ} `eat' with ventive suffix}{}{\cite[133]{frajzyngierGrammarPero1989}} \\
\gll n-káw-tù	mùngbúddínà n-wáat-tù	ɓwá'irò	n	tì	mùngbúddínà-i n-áɗɗ-ínà \\
and-gather-\gl{vent}	elders and-go/come-\gl{vent}	divide.\gl{pl}	\gl{ben}	\gl{prep}	elders-\gl{def} and-eat.\gl{pl}-\gl{compl}.\gl{vent} \\
\glt `(He) will gather the elders and (when) they have gathered he will distribute to the elders (all the meat) and they will eat it.'
\z

\citet[250]{frajzyngierGrammarPero1989} also describes a directional serial verb construction in which the motion verb \textit{ẃaatò} follows a manner of motion verb, as in (\ref{ex2:pero8}) and (\ref{ex2:pero9}).  It is unknown why the manner-of-motion predicates used in Frajzyngier’s examples are preceded by a verb glossed `make'. 

\ea\label{ex2:pero8}
\langinfo{Pero directional multiverb construction}{}{\cite[251]{frajzyngierGrammarPero1989}} \\
\gll yù	káyò	wáatò	pòk-kómbò	 \\
make	swim	go/come	edge-shore \\
\glt `Swim to the other shore.'

\ex\label{ex2:pero9}
\langinfo{Pero directional multiverb construction}{}{\cite[251]{frajzyngierGrammarPero1989}} \\
\gll nì-yí-nà	káyò	wáat-tù	píccè \\
1\gl{sg}-make-\gl{compl}.\gl{vent}	swim	go/come-\gl{vent}	here	 \\
\glt `I swam from there here.'
\z 

\subparagraph{Caused accompanied motion (CAM)}



The ventive form of the verb \textit{íp} `catch' expresses directed caused accompanied motion (CAM), as in (\ref{ex2:pero16}). \citet[96]{frajzyngierGrammarPero1989} suggests that the CAM translation may be an implicature: ``Presumably in this set of sentences the objects have been subsequently brought to the place of speech…''

\ea\label{ex2:pero16}
\langinfo{Pero verb \textit{íp} `catch' with ventive suffix}{}{\cite[99]{frajzyngierGrammarPero1989}} \\
\gll nì-mà-wée-kò	cígí-nì	nì-tà-íp-tù \\
1\gl{sg}-\gl{cond}-see-\gl{compl}	body-3\gl{m}	1\gl{sg}-\gl{fut}-catch-\gl{vent} \\
\glt `If I see him I will catch him and bring him here.'
\z 

The Pero data also shows three possible ways of modifying the motion verb \textit{ẃaatò} `go/come' to express directed CAM: adding an associative prepositional phrase, as in (\ref{ex2:pero17}), using a ``causative'' suffix \textit{-in}, as in (\ref{ex2:pero18}), and the use of a dative (benefactive) pronominal suffix, as in (\ref{ex2:pero19}).

\ea\label{ex2:pero17}
\langinfo{Pero verb \textit{ẃaatò} `go/come' with associative prepositional phrase}{}{\cite[174]{frajzyngierGrammarPero1989}} \\
\gll nì-tá-ẃaat-tù kà cándè	\\
l\gl{sg}-\gl{fut}-go/come-\gl{vent} \gl{asoc} yam		\\
\glt `I will bring yam.'

\ex\label{ex2:pero18}
\langinfo{Pero verb \textit{ẃaatò} `go/come' with causative suffix}{}{\cite[266]{frajzyngierGrammarPero1989}} \\ 
\gll mà-wáat-nà-n mínà wát-tù píppúnò cígí-cù ɗóè  \\
\gl{cond}-go/come-\gl{vent}.\gl{compl}-\gl{tr} home go/come-\gl{vent} wash-\gl{pl} body-\gl{pl} all \\
\glt `When they bring them home they wash themselves completely.'

\ex\label{ex2:pero19}
\langinfo{Pero verb \textit{ẃaatò} `go/come' with dative pronoun}{}{\cite[111]{frajzyngierGrammarPero1989}} \\ 
\gll cì-tà-wát-tù-éenò	\\
2\gl{f}-\gl{fut}-go/come-\gl{vent}-1\gl{sg}(.\gl{dat})	\\	
\glt `You should bring for me.'
\z 	

\newpage
The ``causative'' suffix \textit{-n} is a valency-increasing suffix with several related functions. The functions ascribed to this suffix include causation, benefactive, transitivizing and anaphora \citep[169--176]{frajzyngierGrammarPero1989}. \citet[173]{frajzyngierGrammarPero1989} only gives one unambiguous example of the causative function of \textit{-ni} and notes that the causative function does not appear with transitive verbs. In expressing directed CAM, the \textit{-n} suffix often co-occurs with an associative prepositional phrase, as in (\ref{ex2:pero20}).\footnote{\citet[174]{frajzyngierGrammarPero1989} claims that the ``causative'' suffix indicates that the subject does not accompany the item moving, in which case these should be examples of caused motion rather than caused accompanied motion. However, the translations consistently indicate a caused accompanied motion interpretation, contradicting the proposed analysis.} 

\ea\label{ex2:pero20}
\langinfo{Pero verb \textit{ẃaatò} `go/come' with causative suffix and associative prepositional phrase}{}{\cite[104]{frajzyngierGrammarPero1989}} \\
\gll íkkà wát-tù-n kà cíinà	\\
\gl{prog} go/come-\gl{vent}-\gl{caus} \gl{prep} food \\		
\glt `She is bringing food.'
\z 

\subparagraph{Buy-sell verbs}



The verb \textit{pílù} is glossed `buy' and when it takes the valency-increasing suffix \textit{-n} it is glossed `sell' \citep[175]{frajzyngierGrammarPero1989}. The suffix was earlier labeled a ``causative'' suffix but a more recent analysis labels it an ``additional argument marker'' \citep[84]{frajzyngierLocativePredicationsChadic2024}.

The ventive suffix \textit{-tù} appears in one example,  on the verb \textit{pílù} `buy' with no apparent change in meaning, as shown in (\ref{ex2:pero21}). However, \citet[89]{frajzyngierLocativePredicationsChadic2024} claims that in this example ``the ventive implies movement of the thing bought to the subject.''\footnote{See \citet{lovestrandreview2025} for a review of \citet[]{frajzyngierLocativePredicationsChadic2024}.}

\ea\label{ex2:pero21}
\langinfo{Pero verb \textit{pílù} `buy' with ventive suffix}{}{\cite[9.]{frajzyngierGrammarPero1989}} \\
\gll à píl-tù cándè-m	\\
\gl{neg} buy-\gl{vent} yam-\gl{neg} \\		
\glt `He didn't buy a yam.'
\z 

\paragraph*{Tangale (nucl1696)}  \label{tang}

Information on Tangale grammar is from a lexicon with a grammar sketch \citep{jungraithmayrDictionaryTangaleLanguage1991}. 

\newpage
\subparagraph{Directional verbs}~



\begin{table}[H]
	\caption{Tangale directional verbs} 
	\label{tab:tangV}
	\begin{tabular}{lll} % add l for every additional column or remove as necessary
		\lsptoprule
		DIR & verb & gloss   \\ %table header
		\midrule
		MOTION & \textit{war} & `go' \\
		INTO& \textit{kẹn-} & `enter' \\
		OUT & \textit{pọd-} & `go out' \\
		UP & \textit{ádị} & `climb, go up' \\
		DOWN & \textit{yẹk-} & `go down' \\
		\lspbottomrule
	\end{tabular}
\end{table}

\subparagraph{Directional extensions}~



\begin{table}[H]
	\caption{Tangale directional morphology}
	\label{tab:tang}
	\begin{tabular}{llll} % add l for every additional column or remove as necessary
		\lsptoprule
		forms & source gloss & direction  & other functions \\ %table header
		\midrule
		\textit{-tu, -na}  &  altrilocal-ventive  &  \textsc{vent}  &     \\
		\lspbottomrule
	\end{tabular}
\end{table} 

\citet[46--47]{jungraithmayrDictionaryTangaleLanguage1991} describes ``the altrilocal-ventive or distance stem'' in Tangale. The ventive forms of several intransitive motion verbs illustrate the directional meaning of this suffix, as in (\ref{ex2:tang1}) and (\ref{ex2:tang2}). 

\ea\label{ex2:tang1}
\langinfo{Tangale verb \textit{yẹk} `go down' with ventive suffix}{}{\cite[47]{jungraithmayrDictionaryTangaleLanguage1991}} \\
\gll n yẹk-tụ́-ngọ	\\
1\gl{sg} descend-\gl{vent}-1\gl{sg}	\\
\glt `I came down.'

\ex\label{ex2:tang2}
\langinfo{Tangale verb \textit{pọ́} `go out' with ventive suffix}{}{\cite[47]{jungraithmayrDictionaryTangaleLanguage1991}} \\
\gll pọ́-tụ́-kọ	\\
exit-\gl{vent}-??	\\
\glt `Come out!'
\z 

There are a number of non-motion verbs shown with a ventive suffix, but no translations are given so it is unclear how the ventive forms of these verbs are interpreted in actual use. 

\subparagraph{Caused accompanied motion (CAM)}



The verb \textit{adị} `pass, take, carry' can express caused accompanied motion but no examples were found of directed CAM.

\subparagraph{Buy-sell verbs}



The verb \textit{payị} can have the meanings `trade, buy, sell' \citep[129]{jungraithmayrDictionaryTangaleLanguage1991}. The notion `sell' can also be expressed by the verbs \textit{sọgị} `spend, get rid of, lose, sell' and \textit{wayị} `let go, release, sell' \citep[146, 161]{jungraithmayrDictionaryTangaleLanguage1991}.

\subsection{West Chadic A3 (3/14)}

\citet[308]{schuhChadicCornucopia2017} notes that there is no evidence of directional verbal morphology in West Chadic A3, West Chadic B2, or West Chadic B3 (Schuh’s West Chadic C). \citet{schuhChadicCornucopia2017} discusses four of the 14 West Chadic A3 languages: Goemai (goem1240), Mushere (cakf1236), Mupun (mwag1236) and Ngas (ngas1240). Only three of those are discussed here, as the available analysis of Mushere grammar is limited. In regard to the other ten West Chadic A3 languages, the sparse documentation available does not give evidence of directional verbal morphology, so Schuh's claim is likely correct. For reference, the other ten A3 languages are: Belneng (beln1234), Jakkatoe (jort1240), Koenoem (koen1239), Miship (mish1244), Montol (mont1280), Nteng (nucl1698), Pan (kofy1242), Pyapun (pyap1239), Tal (tall1250) and Yiwom (yiwo1237). 

\paragraph*{Goemai (goem1240)} \label{goem}

Information on the grammar of Goemai is from a reference grammar \citep{hellwigGrammarGoemai2011}.

\subparagraph{Directional verbs}~


\begin{table}[H]
	\caption{Goemai directional verbs} 
	\label{tab:goemV}
	\begin{tabular}{lll} % add l for every additional column or remove as necessary
		\lsptoprule
		DIR & verb & gloss   \\ %table header
		\midrule
		ITIVE & \textit{mu̠ààn} & `go' \\
		VENT & \textit{dóe} & `come' \\
		INTO& \textit{rú} & `enter' \\
		OUT & \textit{p'ét} & `exit' \\
		UP & \textit{hààn} & `climb, ascend' \\
		DOWN & \textit{sàm} & `descend' \\
		\lspbottomrule
	\end{tabular}
\end{table}

\subparagraph{Directional extensions}~

There are no directional extensions in Goemai \citep[10, 172]{hellwigGrammarGoemai2011}.

\subparagraph{Caused accompanied motion (CAM)}



Directed caused accompanied motion can be expressed by the use of a directional verb with a comitative prepositional phrase, as in (\ref{ex2:goem1}). 

\ea\label{ex2:goem1}
\langinfo{Goemai verb \textit{bà} `return' with comitative prepositional phrase}{}{\cite[496]{hellwigGrammarGoemai2011}} \\
\gll D'à góe=kát s'óe, bà ǹ-ní ǹ-lú. \\
\gl{cond} 2\gl{sg}.\gl{m}.\gl{sbj}=find food return(\gl{sg}) with-3\gl{sg} \gl{loc}-settlement \\
\glt `If you find food, you bring it back to the village.'
\z 

\subparagraph{Buy-sell verbs}



\citet[214]{hellwigGrammarGoemai2011} analyzes the disambiguation of the meanings of the verb \textit{s'éét} `buy/sell'. In (\ref{ex2:goem2}), the `sell' interpretation is required by the number agreement between the verb and the comitative prepositional phrase. This agreement with a non-core argument can only occur in the ``causative construction'' which gives rise to a `sell' interpretation rather than `buy'. 

\ea\label{ex2:goem2}
\langinfo{Goemai verb \textit{s'eet} `buy/sell' with comitative prepositional phrase}{}{\cite[217]{hellwigGrammarGoemai2011}} \\
\gll mòe=s'èèt m̠ép góe lá=kè \\
l\gl{pl}.\gl{sbj}=buy/sell(\gl{sg}) 3\gl{pl}.\gl{obj} with \gl{dem}(\gl{sg}).\gl{gen}=chicken \\
\glt `We sold them the little chicken.'
\z 

However, strictly speaking, the presence of a comitative prepositional phrase on its own does not force the `sell' interpretation, as the comitative phrase could (in the right context) still be interpreted as an adjunct, rather than as a non-core argument. This interpretation is the second translation of (\ref{ex2:goem3}).

\ea\label{ex2:goem3}
\langinfo{Goemai verb \textit{s'eet} `buy/sell' with comitative prepositional phrase}{}{\cite[216]{hellwigGrammarGoemai2011}} \\
\gll  ní s'éét hèn góe shím \\
3\gl{sg}.\gl{sbj} buy/sell(\gl{sg}) 1\gl{sg}.\gl{obj} with yam \\
\glt `He sold me yam.' or `He bought me (over) with yam.'
\z 

\paragraph*{Mupun (Mwaghavul) (mwag1236)} \label{mupu}

\citet{frajzyngierGrammarMupun1993} is a reference grammar of a dialect of Mwaghavul called Mupun.

\newpage
\subparagraph{Directional verbs}~



\begin{table}[H]
	\caption{Mupun directional verbs} 
	\label{tab:mupuV}
	\begin{tabular}{lll} % add l for every additional column or remove as necessary
		\lsptoprule
		DIR & verb & gloss   \\ %table header
		\midrule
		ITIVE & \textit{dəm} & `go' \\
		VENT & \textit{ji} & `come' \\
		UP & \textit{ka} & `ascend' \\
		DOWN & \textit{siam} & `descend' \\
		INTO& \textit{ɗel} & `enter' \\
		OUT & \textit{pūt} & `go out' \\
		\lspbottomrule
	\end{tabular}
\end{table}

\subparagraph{Directional extensions}



There are no directional extensions in Mupun, but both directional meaning and prior associated motion can be expressed through serial verb constructions. There are examples of directional serial verb constructions which show a directional verb following another motion verb. The directional verb can indicate itive motion, as in (\ref{ex2:mupu6}), ventive motion, as in (\ref{ex2:mupu7}), or downward motion, as in (\ref{ex2:mupu8}). 

\ea\label{ex2:mupu6}
\langinfo{Mupun itive directional serial verb construction}{}{\cite[241]{frajzyngierGrammarMupun1993}} \\
\gll  an	mbə	wa	dəm	n-tul \\
1\gl{sg}	\gl{fut}	return	go	\gl{prep}-home \\
\glt `I will return home.'

\ex\label{ex2:mupu7}
\langinfo{Mupun ventive directional serial verb construction}{}{\cite[241]{frajzyngierGrammarMupun1993}} \\
\gll  wa	n-baa	n-ji	n-tul \\
return	1\gl{sg}-return	1\gl{sg}-come	\gl{prep}-home \\
\glt `I returned home. (If the first person subject is already at the place.)'

\ex\label{ex2:mupu8}
\langinfo{Mupun downward directional serial verb construction}{}{\cite[229]{frajzyngierGrammarMupun1993}} \\
\gll  wa	mu	siam	n-tulu \\
return	1\gl{pl}	descend	\gl{prep}-home \\
\glt `We went down home.'
\z 

\subparagraph{Caused accompanied motion (CAM)}



It appears that caused accompanied motion in Mupun is expressed using a motion verb and a prepositional phrase, as in (\ref{ex2:mupu13}).

\ea\label{ex2:mupu13}
\langinfo{Mupun caused accompanied motion}{}{\cite[277]{frajzyngierGrammarMupun1993}} \\
\gll wu ji a kə daam məndoŋ ɓejee/dak \\
3M come \gl{cop} \gl{prep} bag one ?? \\
\glt `He brought only one bag.'
\z 

\subparagraph{Buy-sell verbs}



The verb \textit{seet} is glossed `buy' and `sell'. There is no evidence of a grammaticalized means of distinguishing the two interpretations.

\paragraph*{Ngas (ngas1240)} \label{ngas} Information on Ngas grammar is primarily from a lexicon and text collection with a short grammar sketch \citep{jungraithmayrNgasLanguageShik2016}. 

\subparagraph{Directional verbs}~



\begin{table}[H]
	\caption{Ngas directional verbs} 
	\label{tab:ngasV}
	\begin{tabular}{lll} % add l for every additional column or remove as necessary
		\lsptoprule
		DIR & verb & gloss   \\ %table header
		\midrule
		ITIVE & \textit{dəm, met} & `go' \\
		VENT & \textit{ji} & `come' \\
		DOWN & \textit{tum, dam} & `descend' \\
		INTO& \textit{sit} & `enter' \\
		OUT & \textit{put} & `go out' \\
		\lspbottomrule
	\end{tabular}
\end{table}

Upward motion is expressed using an adverbial expression \textit{kə ɠyéeŋ} \citep[104]{jungraithmayrNgasLanguageShik2016}.

\subparagraph{Directional extensions}~



Publications on Ngas do not mention any ventive or other directional marker \citep{burquestGrammarAngas1973,jungraithmayrNgasLanguageShik2016}. 

\subparagraph{Caused accompanied motion (CAM)}



Ventive caused accompanied motion (CAM) is expressed by a ventive motion verb with a prepositional phrase of accompaniment: \textit{jə kə} `come with, bring' \citep[182]{jungraithmayrNgasLanguageShik2016}. The Ngas lexicon includes the verb \textit{ɓam} `to take away, snatch, seize; rescue, help, save' \citep[162]{jungraithmayrNgasLanguageShik2016} which could conceivably be construed as an itive CAM verb. However, since most of the words used to describe its meaning relate to the idea of possession rather than translational motion, it has not been included as a CAM verb.

\subparagraph{Buy-sell verbs}



In Ngas, the verb \textit{sit} [síit] means `buy' and \textit{ser} [séer] means `sell' \citep[216--217]{jungraithmayrNgasLanguageShik2016}. 

\subsection{West Chadic A4 (4/8)}

Four of the eight West Chadic A4 languages are discussed here. For Tambas (tamb1267), Mangar (mang1417), Duhwa (duhw1236) or Mindat (mund1334), no relevant descriptions were found.

\paragraph*{Fyer (fyer1241)}

Information on the grammar of Fyer is from a grammar sketch included as a chapter in a book on several West Chadic A4 languages \citep{jungraithmayrRonSprachenTaschadohamitischeStudien1970}.

\subparagraph{Directional verbs}~



\begin{table}[H]
	\caption{Fyer directional verbs} 
	\label{tab:fyerV}
	\begin{tabular}{lll} % add l for every additional column or remove as necessary
		\lsptoprule
		DIR & verb & gloss   \\ %table header
		\midrule
		ITIVE & \textit{wu, gan} & `\textit{gehen} [go]' \\
		VENT & \textit{'èl} & `\textit{kommen} [come]' \\
        INTO& \textit{ton} & `\textit{eintreten}, \textit{betreten} [go in, enter]' \\		OUT &  \textit{saŋ} & `\textit{ausgehen} [go out]' \\
		UP & \textit{ɗoo} & `\textit{klettern}, \textit{aufsteigen} [climb, rise up]' \\
		DOWN & \textit{ɗwaŋ} & `\textit{absteigen} [go down]' \\
		\lspbottomrule
	\end{tabular}
\end{table}

\subparagraph{Directional extensions}


No description of motion-related verbal morphology was found in the sketch of Fyer grammar by \citet[15--96]{jungraithmayrRonSprachenTaschadohamitischeStudien1970}.

\subparagraph{Caused accompanied motion (CAM)}



Insufficient information found. 

\subparagraph{Buy-sell verbs}



Fyer has a verb \textit{gon} `buy and sell, trade' [German: `\textit{kaufen und verkaufen, handeln}'] but no explanation is given for how the notions BUY and SELL might be differentiated in the language \citep[85]{jungraithmayrRonSprachenTaschadohamitischeStudien1970}. 

\paragraph*{Kulere (kule1247)}

Information on the grammar of Kulere is from a grammar sketch included as a chapter in a book on several West Chadic A4 languages \citep{jungraithmayrRonSprachenTaschadohamitischeStudien1970}.

\subparagraph{Directional verbs}~



\begin{table}[H]
	\caption{Kulere directional verbs} 
	\label{tab:kuleV}
	\begin{tabular}{lll} % add l for every additional column or remove as necessary
		\lsptoprule
		DIR & verb & gloss   \\ %table header
		\midrule
		ITIVE & \textit{fa} & `\textit{gehen} [go]' \\
		VENT & \textit{bo} & `\textit{kommen} [come]' \\
		INTO& \textit{ra} & `\textit{eintreten} [go in]' \\
		OUT & \textit{laŋ} & `\textit{ausgehen} [go out]' \\
		UP & \textit{regy} & `\textit{klettern}, \textit{aufsteigen} [climb, rise up]' \\
		DOWN & \textit{ɗor-o} & `\textit{absteigen} [go down]' \\
		\lspbottomrule
	\end{tabular}
\end{table}

\subparagraph{Directional extensions}



No description of motion-related verbal morphology was found in the sketch of Kulere grammar by \citet[295--360]{jungraithmayrRonSprachenTaschadohamitischeStudien1970}.

\subparagraph{Caused accompanied motion (CAM)}



Insufficient information found. 

\subparagraph{Buy-sell verbs}



The verb \textit{dyef} is glossed `buy, trade' [German: `\textit{kaufen, handeln}'] and its entry in the lexicon notes a derived form \textit{dyef ɗes} `sell' \citep[351]{jungraithmayrRonSprachenTaschadohamitischeStudien1970}, however, it is unclear what type of morpheme \textit{ɗes} is.

\paragraph*{Ron-Daffo (ronn1241)}

Information on the grammar of Ron-Daffo is from a grammar sketch included as a chapter in a book on several West Chadic A4 languages \citep{jungraithmayrRonSprachenTaschadohamitischeStudien1970} and from a reference grammar \citep[43]{seibertRonDaffoJosPlateau1998}.

\subparagraph{Directional verbs}~



\begin{table}[H]
	\caption{Ron-Daffo directional verbs} 
	\label{tab:rondV}
	\begin{tabular}{lll} % add l for every additional column or remove as necessary
		\lsptoprule
		DIR & verb & gloss   \\ %table header
		\midrule
		ITIVE & \textit{yû} & `\textit{gehen} [go]'  \\
		VENT & \textit{yes} & `\textit{kommen} [come]' \\
		UP & \textit{ɗu, lâŋ} & `(\textit{be})\textit{steigen} [ascend]' \\
		DOWN & \textit{ɗor} & `\textit{herabsteigen} [descend]' \\
		\lspbottomrule
	\end{tabular}
\end{table}

\subparagraph{Directional extensions}



\citet[119]{jungraithmayrRonSprachenTaschadohamitischeStudien1970} and \citet[43]{seibertRonDaffoJosPlateau1998} state that there are no productive directional verbal extensions in Ron-Daffo, but there are three cases where directional motion predicates are formed from the verb \textit{lâŋ} `go up, get up' combined with a particle or adposition: \textit{láŋ na} `come in' [German: `\textit{hereinkommen}'], \textit{lâŋ ti} `enter', \textit{lâŋ lá} `go out' [German: `\textit{ausgehen}']. This is not a productive process and the semantics of the composed predicates are not transparently compositional. 

\subparagraph{Caused accompanied motion (CAM)}



Caused accompanied motion can be expressed by the verb \textit{kùl} `bring'. Directed CAM can be expressed by the use of a directional verb with an adposition \textit{ti} `with' as in: \textit{yû ti} `bring here' [German: `\textit{hinbringen}'] and \textit{yes ti} `bring along' [German: `\textit{mitbringen}'] \citep[43]{seibertRonDaffoJosPlateau1998}. 

\subparagraph{Buy-sell verbs}



The verb \textit{gôn} is glossed `buy'. With the particle \textit{lá} (equivalent to the Hausa totality form ``Grade 4'', but also expressing itive direction), the predicate \textit{gôn lá} means `sell'. 

\paragraph*{Ron-Scha (shaa1247)} \label{shaa}

Information on the grammar of Ron-Scha is from a grammar sketch included as a chapter in a book on several West Chadic A4 languages \citep{jungraithmayrRonSprachenTaschadohamitischeStudien1970}.

\subparagraph{Directional verbs}~



\begin{table}[H]
	\caption{Ron-Scha directional verbs} 
	\label{tab:shaaV}
	\begin{tabular}{lll} % add l for every additional column or remove as necessary
		\lsptoprule
		DIR & verb & gloss   \\ %table header
		\midrule
		ITIVE & \textit{du, fay} & `\textit{gehen} [go]' \\
		VENT & \textit{bol} & `\textit{kommen} [come]' \\
		UP & \textit{'âg} & `\textit{hinaufsteigen} [ascend]' \\
		DOWN & \textit{dôh} & `\textit{absteigen} [descend]' \\
		INTO/OUT & \textit{ḥoh} & `\textit{eintreten} [enter]' \\
		\lspbottomrule
	\end{tabular}
\end{table}

Outward motion is expressed as \textit{ḥóh kweŋ}. Since the verb \textit{ḥóh} `enter' can express outward motion when combined with the adverbial \textit{kweŋ} `out' [German: `\textit{hinaus}'], it appears to be a general boundary-crossing verb. It is possible that the expression \textit{ḥóh kweŋ} `exit' should be treated as a lexicalized form, but here it is assumed to be evidence that in the right semantic context the verb \textit{ḥóh} can be shown to be ambiguous between inward and outward motion.

\subparagraph{Directional extensions}~



\begin{table}[H]
	\caption{Ron-Scha directional morphology}
	\label{tab:shaa}
	\begin{tabular}{llll} % add l for every additional column or remove as necessary
		\lsptoprule
		forms & gloss & direction  & other functions \\ %table header
		\midrule
		\textit{-ó}  & \textsc{vent} &  \textsc{vent} &  \textsc{loc}  \\
		\lspbottomrule
	\end{tabular}
\end{table} 

Ron-Scha (or Sha) is described by \citet[268–269]{jungraithmayrRonSprachenTaschadohamitischeStudien1970} as having a ventive suffix \textit{-ó} comparable to the Hausa ventive form \citep[cf.][308]{schuhChadicCornucopia2017}. Examples are given of ventive forms of intransitive and transitive motion verbs, as in (\ref{ex2:shaa1}) and (\ref{ex2:shaa2}). 

\ea\label{ex2:shaa1}
\langinfo{Sha verb \textit{doḥ} `descend' with ventive suffix}{}{\cite[268]{jungraithmayrRonSprachenTaschadohamitischeStudien1970}} \\
\gll doḥ-ó	\\ 
descend-\gl{vent} \\
\glt `descend here' [German: `\textit{her-absteigen}']

\ex\label{ex2:shaa2}
\langinfo{Sha verb \textit{də̀ŋ} `pull' with ventive suffix}{}{\cite[268]{jungraithmayrRonSprachenTaschadohamitischeStudien1970}} \\
\gll də̀ŋ-ó \\
pull-\gl{vent} \\
\glt `drag here' [German: `\textit{her-ziehen}']
\z 

The examples given of the \textit{-ó} suffix are all with motion verbs with the exception of three verb roots that are not attested without the ventive suffix. This is taken as lexicalization of the ventive form, as in (\ref{ex2:shaa3}) and (\ref{ex2:shaa4}). 

\ea\label{ex2:shaa3}
\langinfo{Sha verb \textit{cə̀l} `catch' with (obligatory) ventive suffix}{}{\cite[269]{jungraithmayrRonSprachenTaschadohamitischeStudien1970}} \\
\gll cə̀l-ó \\
??-\gl{vent} \\
\glt `catch here' [German: `{``}\textit{her}-''\textit{fangen}']

\ex\label{ex2:shaa4}
\langinfo{Sha verb \textit{wo'} `wake up' with (obligatory) ventive suffix}{}{\cite[269]{jungraithmayrRonSprachenTaschadohamitischeStudien1970}} \\
\gll wo'-ó(h) \\
??-\gl{vent} \\
\glt `wake up' [German: `{``}\textit{her}''-\textit{aufwecken}']
\z 

\subparagraph{Caused accompanied motion (CAM)}



There is no information available about caused accompanied motion in Sha. 

\subparagraph{Buy-sell verbs}



Both buying and selling are expressed by the same verb, \textit{hyù} `\textit{kaufen} [buy], \textit{verkaufen} [sell]' \citep[292, 294]{jungraithmayrRonSprachenTaschadohamitischeStudien1970}, with no data given on how the two interpretations might be disambiguated. 

\section{West Chadic B (7/38)}

\subsection{West Chadic B1 (2/5)}

Two West Chadic B1 languages are included here: Bade and Ngizim. No relevant descriptions were found for Duwai (duwa1244), Auyokawa (auyo1240) or Teshenawa (tesh1239). The latter two are considered dormant. 

\paragraph*{Bade (bade1248)} \label{bade} Information on Bade grammar is primarily from a chapter in an edited volume giving an overview of the morphology \citep[]{schuhBadeMorphology2007}. 

\subparagraph{Directional verbs}~



\begin{table}[H]
	\caption{Bade directional verbs} 
	\label{tab:badeV}
	\begin{tabular}{lll} % add l for every additional column or remove as necessary
		\lsptoprule
		DIR & verb & gloss   \\ %table header
		\midrule
		ITIVE & \textit{ju} & `go' \\
		VENT & \textit{dàawau} & `come' \\
		INTO& \textit{ə̀fku} & `enter' \\
		OUT & \textit{və̀ru} & `go out' \\
		\lspbottomrule
	\end{tabular}
\end{table}

\subparagraph{Directional extensions}~



\begin{table}[H]
	\caption{Bade directional morphology}
	\label{tab:bade}
	\begin{tabular}{llll} % add l for every additional column or remove as necessary
		\lsptoprule
		forms & source gloss & direction  & other functions \\ %table header
		\midrule
		\textit{-àawo, -ìina}  &  ventive  &   \textsc{vent} &   \textsc{subs.vent}  \\
		\lspbottomrule
	\end{tabular}
\end{table} 

\citet[620–621]{schuhBadeMorphology2007} describes ventive forms in Bade giving just a few examples with translations (but not interlinear glossing) such as \textit{gə və̀rā̀wo} `you came out (speaker is outside)' and \textit{nə̂-bzā̀wo} `I left (it back there and came here)'. (The macron denotes a long vowel.) The former example appears to be a directional interpretation with a motion verb, and the latter appears to be a subsequent motion interpretation with a non-motion verb. \citet[307]{schuhChadicCornucopia2017} also includes two examples of ventive forms: the Perfective form of the verb `go out, exit' with a ventive suffix, \textit{və̀r-àawo}, and the ``second subjunctive'' form of the same verb with a ventive suffix, \textit{də̀ və̀r-ìina}. \citet[14]{ziegelmeyerVerbalSystemGashua2014} identifies a similar form in the Gashua dialect of Bade which he describes as a suffix: \textit{-àawo}. 

\subparagraph{Caused accompanied motion (CAM)}



It appears that motion verbs can combine with a valency-increasing suffix \textit{-du} to express directed caused accompanied motion. For example, the verb \textit{ju} `go' becomes \textit{jə-dù} `take, transport, carry' \citep[618]{schuhBadeMorphology2007}, and the verb \textit{vǝ̀lu} `go out' becomes \textit{vǝ̀lǝ-dù} `bring out something' \citep[13]{ziegelmeyerVerbalSystemGashua2014}.

\subparagraph{Buy-sell verbs}



Bade has two separate verbs for buying and selling: \textit{màsu} `buy' and \textit{ɗə̀bdiu} `sell'. 

\paragraph*{Ngizim (ngiz1242)} \label{ngiz} Information on Ngizim grammar is from the PhD dissertation of \citet{schuhAspectsNgizimSyntax1972}.

\subparagraph{Directional verbs}~



\begin{table}[H]
	\caption{Ngizim directional verbs} 
	\label{tab:ngizV}
	\begin{tabular}{lll} % add l for every additional column or remove as necessary
		\lsptoprule
		DIR & verb & gloss   \\ %table header
		\midrule
		ITIVE & \textit{jú-w} & `go' \\
		VENT & \textit{dee-w, nài} & `come' \\ 
		INTO& \textit{tə̀fi} & `enter' \\
		OUT & \textit{və̀rà} & `go out' \\
		\lspbottomrule
	\end{tabular}
\end{table}

\subparagraph{Directional extensions}~



\begin{table}[H]
	\caption{Ngizim directional morphology}
	\label{tab:ngiz}
	\begin{tabular}{llll} % add l for every additional column or remove as necessary
		\lsptoprule
		forms & source gloss & direction  & other functions \\ %table header
		\midrule
		\textit{-ay, -iina}  &  \textsc{vent}  &   \textsc{vent} &   \textsc{subs.vent}  \\
		\lspbottomrule
	\end{tabular}
\end{table} 

\citet[26]{schuhAspectsNgizimSyntax1972} describes ventive forms in Ngizim as a set of irregular TAM suffixes and possible tone change: \begin{quote}The Ventive is always indicated by replacement of the neutral verb suffix by a special suffix, characterizable as AY, or in some aspects, \textit{-iina}. There are slight variations on AY conditioned by aspect. In addition to this suffix change, all verbs in the Subjunctive and Imperative take low tone, regardless of lexical tone, when the Ventive is added.\end{quote} 

The function of the ventive is broadly described as follows: \begin{quote}
This extension indicates action which takes place in the direction of, or for the benefit of, some person or place of reference (often the speaker or the speaker's location). Sometimes the event itself may have taken place somewhere else, but it ultimately affects the place or person of reference, e.g., `ask' + Ventive = `ask and bring back the answer' \citep[26]{schuhAspectsNgizimSyntax1972}.\end{quote} This definition indicates that not only does the ventive have a directional function, as in (\ref{ex2:ngiz1}) and (\ref{ex2:ngiz2}), but that the subsequent motion interpretation can be found with the ventive form of non-motion verbs such as `buy' in (\ref{ex2:ngiz3}). 

\ea\label{ex2:ngiz1}
\langinfo{Ngizim verb \textit{tə̀f} `enter' with ventive suffix}{}{\cite[26]{schuhAspectsNgizimSyntax1972}} \\
\gll  tə̀f-éew \\
enter-\gl{vent}.\gl{pfv} \\
\glt `He came in.'

\ex\label{ex2:ngiz2}
\langinfo{Ngizim verb \textit{wàn} `send' with ventive suffix}{}{\cite[27]{schuhAspectsNgizimSyntax1972}} \\
\gll  á	wàn-ìiná \\
\gl{aux}	send-\gl{vent}.\gl{imp}.\gl{pl} \\
\glt `Send [\textsc{pl}] (here)!'

\ex\label{ex2:ngiz3}
\langinfo{Ngizim verb \textit{màs} `buy' with ventive suffix}{}{\cite[26]{schuhAspectsNgizimSyntax1972}} \\
\gll jà	màs-ée	máraû \\
1\gl{pl}	buy.\gl{vent}-\gl{pfv}	millet \\
\glt `We bought (and brought) millet.'
\z 

\subparagraph{Caused accompanied motion (CAM)}



Directed caused accompanied motion can be expressed by the ventive form of the verb \textit{jìb} `catch', as in (\ref{ex2:ngiz4}).

\ea\label{ex2:ngiz4}
\langinfo{Ngizim verb \textit{jìb} `catch' with ventive suffix}{}{\cite[91]{schuhAspectsNgizimSyntax1972}} \\
\gll ...ká	jìb-ài	ácî \\
2\gl{sg}	catch-\gl{vent}.\gl{sbjv}	3\gl{sg} \\
\glt `...that you catch (and bring) him.'
\z 

\subparagraph{Buy-sell verbs}



There are two separate verbs for buying and selling: \textit{másə} `buy' and \textit{ɗə́bdə} `sell' \citep[14]{schuhAspectsNgizimSyntax1972}.

\subsection{West Chadic B2 (2/10)}

Grammatical descriptions are available for only two of the West Chadic B2 languages (also called North Bauchi): Pa'a and Miya. The languages not included here are: Ciwogai (ciwo1236), Diri (diri1259), Mburku (mbur1239), Kariya (kari1316), Siri (siri1278), Warji (warj1253), Zumbun (zumb1240), and the dormant language Ajawa (ajaw1236). 

\paragraph*{Pa'a (paaa1242)} \label{paaa}

Information on Pa'a grammar is from a doctoral dissertation \citep{skinnerAspectsPaanciGrammar1979}. 

\subparagraph{Directional verbs}~



\begin{table}[H]
	\caption{Pa'a directional verbs} 
	\label{tab:paaaV}
	\begin{tabular}{lll} % add l for every additional column or remove as necessary
		\lsptoprule
		DIR & verb & gloss   \\ %table header
		\midrule
		ITIVE & \textit{ĥwòcú} & `go' \\
		VENT & \textit{dàvá} & `come' \\ 
		INTO& \textit{zaa} & `enter' \\
		OUT & \textit{mba} & `go out' \\
		DOWN & \textit{ɗʉbu} & `descend' \\
		\lspbottomrule
	\end{tabular}
\end{table}

\subparagraph{Directional extensions}


No discussion of motion-related verbal morphology was found in the analysis of \citet{skinnerAspectsPaanciGrammar1979}.

\subparagraph{Caused accompanied motion (CAM)}



The verb \textit{ciyè} `bring' expresses caused accompanied motion (CAM) \citep[171]{skinnerAspectsPaanciGrammar1979}. Directed CAM can be expressed by a motion verb in a causative form. For example, the causative form of \textit{mba} `go' is \textit{mbei} `take out' and the causative form of \textit{zaa} `enter' is \textit{zei} `put in' \citep[131]{skinnerAspectsPaanciGrammar1979}. 

\subparagraph{Buy-sell verbs}



\citet[189]{skinnerAspectsPaanciGrammar1979} notes that in regard to the verb \textit{kwan} `buy, sell s.t.' ``context determines meaning, no ± caus. difference.'' 

\paragraph*{Miya (miya1266)}

Information on the grammar of Miya is from a reference grammar \citep{schuhGrammarMiya1998}. 

\subparagraph{Directional verbs}~



\begin{table}[H]
	\caption{Miya directional verbs} 
	\label{tab:miyaV}
	\begin{tabular}{lll} % add l for every additional column or remove as necessary
		\lsptoprule
		DIR & verb & gloss   \\ %table header
		\midrule
		ITIVE & \textit{ba} & `go' \\
		VENT & \textit{buw} & `come' \\
		INTO& \textit{z} & `enter' \\
		OUT & \textit{baw} & `go out' \\
		UP & \textit{ghəma} & `climb, mount' \\
		DOWN & \textit{daw} & `get down; lodge' \\
		\lspbottomrule
	\end{tabular}
\end{table}

\subparagraph{Directional extensions}



\citet[169]{schuhGrammarMiya1998} notes that Miya does not have directional verbal morphology: ``Particularly notable is the absence of a \textit{ventive} stem or \textit{Distanzstamm}, such as Hausa Grade VI marked by an \textit{-oo} termination.'' It is worth noting that the final \textit{w} in the stems of some of the motion verbs in Table \ref{tab:miyaV} do suggest a possible fossilized suffix related to motion meaning. 

\subparagraph{Caused accompanied motion (CAM)}



Transitivized forms of motion verbs can expressed directed CAM as in \textit{buw} `come' and \textit{buway} `bring', \textit{ba} `go' and \textit{bay} `take' and \textit{baw} `go out' and \textit{baway} `take out' \citep[178]{schuhGrammarMiya1998}.

\subparagraph{Buy-sell verbs}



The verb roots for the notions BUY and SELL are unrelated in Miya: \textit{kə̀na} `buy' and \textit{mə́tsə̀} `sell' \citep[169]{schuhGrammarMiya1998}.

\subsection{West Chadic B3 (3/23)}
\sloppy
The West Chadic B3 languages, also called South Bauchi, are divided into a West and an East group. The only B3-West language included here is Guruntum. The remaining B3-West languages are Boghom (bogh1241), Chaari (dans1239), Dokshi (lush1256), Kir-Balar (kirb1236), Mangas (mang1416), Jimi (jimi1255), Ju (juuu1243), Tala (tala1295), Tulai (nucl1693) and Zangwal (zang1255). Two B3-East languages are included: Dass and Saya (or Zaar). The other languages of this subgroup are
Bu (zara1252),
Gyaazi (nucl1692),
Buli (buli1260),
Zul (zull1239),
Dir-Nyamzak-Mbarimi (lund1276),
Pesse (polc1243),
Dyarim (dyar1234) and three languages considered dormant: Luri (luri1256), Zari (zari1242) and Zeem (zeem1242). Several other languages in this group are considered moribund. For several languages there are notes available on some features of the grammar, especially those published by Bernard Caron or Roger Blench. These notes do not mention any sign of directional verbal morphology, but were not considered thorough enough descriptions to be included in the current data set.
\fussy

\paragraph*{Guruntum (guru1271)} \label{guru}

Information on Guruntum grammar is from an extended grammar sketch by \citet{harunaGrammaticalOutlineGurdung2003} and a short sketch by \citet{jaggarGuruntumGurdungWest1988}. 

\subparagraph{Directional verbs}~



\begin{table}[H]
	\caption{Guruntum directional verbs} 
	\label{tab:guruV}
	\begin{tabular}{lll} % add l for every additional column or remove as necessary
		\lsptoprule
		DIR & verb & gloss   \\ %table header
		\midrule
		ITIVE & \textit{karmi} & `go' \\
		VENT & \textit{wari} & `come' \\
		INTO& \textit{ndàa} & `enter' \\
		OUT & \textit{tài} & `go out' \\
		DOWN & \textit{shiwi} & `descend' \\
		\lspbottomrule
	\end{tabular}
\end{table}

Note the verb \textit{taà} `climb' could potentially be an upward motion verb but is glossed as a manner-of-motion verb.

\subparagraph{Directional extensions}



\citet{harunaGrammaticalOutlineGurdung2003} and \citet{jaggarGuruntumGurdungWest1988} do not describe any motion-related verbal morphology in Guruntum. Where directional glosses appear, the morphemes used are described as prepositions. 

\subparagraph{Caused accompanied motion (CAM)}



The verbs \textit{yili} `take' and \textit{wuri} `bring' expressed caused accompanied motion, but it is not clear what expressions are available for expressing directed CAM. 

\subparagraph{Buy-sell verbs}



The notion SELL is derived from the verb \textit{yab} `buy'. According to \citet[173]{jaggarGuruntumGurdungWest1988}, ``The verb + preposition \textit{yab}(\textit{i}) \textit{gʷǎi} ‘sell’ (lit: buy away) is the semantic analogue of the Hausa ‘efferential’ (alias ‘causative’) ‘grade 5’.''

\newpage
\paragraph*{Dass (dass1243)} \label{dass}

Information on the grammar of Dass is from a grammar sketch published as a journal article \citep{caronDottAkaZodi2002}. 

\subparagraph{Directional verbs}~



\begin{table}[H]
	\caption{Dass directional verbs} 
	\label{tab:dassV}
	\begin{tabular}{lll} % add l for every additional column or remove as necessary
		\lsptoprule
		DIR & verb & gloss   \\ %table header
		\midrule
		ITIVE & \textit{rə́s} & `go' \\
		VENT & \textit{tʃetti} & `coming' \\
		INTO/OUT & \textit{tá} & `enter, come out' \\
		UP & \textit{tswaa} & `climb' \\
		DOWN & \textit{ʃi} & `get down' \\
		\lspbottomrule
	\end{tabular}
\end{table}

Since the causative/transitive form of \textit{tswaa} `climb' is \textit{tswár} `take up', it is assumed that the verb is an upward motion verb despite being glossed as a manner-of-motion verb \citep[3]{caronDottAkaZodi2002}. 

\subparagraph{Directional extensions}



\citet[3]{caronDottAkaZodi2002} mentions that no verbal extensions were found in Dass apart from a series of causative/transitivized forms. 

\subparagraph{Caused accompanied motion (CAM)}



In addition the CAM verb \textit{sə́ndər} `bring', \citet[3]{caronDottAkaZodi2002} describes the use of the causative form of motion verbs to express directed CAM, such as \textit{ta} `come out' and \textit{tár} `get out', \textit{ʃi} `get down' and \textit{ʃír} `take down', and \textit{tswaa} `climb' and \textit{tswár} `take up'. 

\subparagraph{Buy-sell verbs}



There are two separate verbs for these notions in Dass: \textit{yep} `buy' and \textit{mándər} `sell'. 

\paragraph*{Zaar (Saya) (saya1246)}

Information on the grammar of Zaar is from an unpublished grammar sketch available as an online manuscript \citep{caronZaarGrammaticalSketch2011}. 

\newpage
\subparagraph{Directional verbs}~



\begin{table}[H]
	\caption{Zaar directional verbs} 
	\label{tab:zaarV}
	\begin{tabular}{lll} % add l for every additional column or remove as necessary
		\lsptoprule
		DIR & verb & gloss   \\ %table header
		\midrule
		ITIVE & \textit{slə} & `go' \\
		VENT & \textit{man} & `come' \\
		INTO/DOWN & \textit{shiː} & `enter' \\
		OUT &\textit{pòtè} &`exit' \\	
		OUT &\textit{nyol} &`go out' \\
		DOWN & \textit{sapká} & `get down' \\
		\lspbottomrule
	\end{tabular}
\end{table}

\subparagraph{Directional extensions}



\citet{caronZaarGrammaticalSketch2011} does not describe any motion-related verbal morphology in Zaar.

\subparagraph{Caused accompanied motion (CAM)}


 
Directed CAM may be expressed by the causative forms of motion verbs, as in \textit{slə} `go' in the causative form \textit{slə́ːr} `take to', \textit{shiː} `enter, get down' in the causative form \textit{shíːr} `take sth. down' or  \textit{nyol} `go out' in the causative form \textit{nyolár} `bring out' \citep[10]{caronZaarGrammaticalSketch2011}.

\subparagraph{Buy-sell verbs}



The notion `sell' is derived from the verb \textit{ɗiːp} `buy' in its causative form \textit{ɗiːbár} `sell' \citep[10]{caronZaarGrammaticalSketch2011}.

